\chapter{辛拓扑中的应用}
\label{ch:applications}

只有借助\personFukaya{}在 (Fu93) 中引入、并在 (FOOO09) 中得到完整发展的
$A_\infty$ 工具，拉格朗日\personFloer{}同调理论的全部威力才能发挥出来。当圆盘气泡的
结构不像以下两种特殊情形那样简单时，处理一般拉格朗日子流形就需要这一理论：
\begin{enumerate}[label=(\arabic*)]
\item （非紧）正合辛流形中的正合拉格朗日子流形；
\item （单调）辛流形中的单调拉格朗日子流形。
\end{enumerate}
由于 $A_\infty$ 工具超出了本书范围，我们将集中讨论这两种情形，并用它们说明\personFloer{}同调在辛拓扑研究中的用法。

不过，即使\personFloer{}同调 $HF(L_0,L_1;M)$ 已经定义，只要超出正合情形以及
$L_1=\phi(L_0)$（其中 $\phi$ 是哈密顿微分同胚）的情形，要显式计算这个同调仍是
一项极为困难的任务。

定理 16.4.10 或它的类似版本，是把\personFloer{}同调用于辛拓扑研究的基本出发点。迄今
多数应用都与构造不可位移拉格朗日子流形，或研究可位移拉格朗日子流形的辛拓扑
有关，例如\personHofer{}位移能与马斯洛夫类障碍。在这些研究中，分析
$\phi\to\mathrm{id}$ 时\personFloer{}模空间的结构也至关重要；更确切地说，是分析
$\phi(L)\to L$ 的绝热极限，而这会产生\personFloer{}流线的厚--薄分解（及相应的\termThickThinDichotomy{}）。

对于希望进一步了解上述两种情形之外的拉格朗日\personFloer{}理论，以及它在镜像对称和
辛拓扑中的应用的读者，我们推荐 (CO06)、(FOOO10b)--(FOOO13)。

\section{邻近拉格朗日对子：\termThickThinDichotomy{}}
\label{sec:nearby-lagrangian-thick-thin}

本节研究当 $\phi\to\mathrm{id}$（按 $C^1$ 拓扑）时
$\mathcal M(L,\phi(L);p,q)$ 的一些收敛结果。我们还将说明，这种退化的研究如何
产生 (Oh96b) 中引入的一个谱序列。这个形式的谱序列后来在 (Bu10)、(BCo09)、
(D09) 中得到进一步研究。(FOOO09) 给出了这类谱序列更一般的版本；在附加\emph{无障碍}
假设下，它也能处理非单调拉格朗日子流形。

下面给出研究这一退化的精确设定。设 $L$ 是给定的紧拉格朗日子流形。以下概念由
(Oh05d, Spa08) 引入。

\stfhIndex{I073}{Tunmeihamiduntonghen}{吞没哈密顿同痕（engulfed Hamiltonian isotopy）}
\begin{definition}
\label{def:v-engulfed-isotopy}
设 $V$ 是 $L$ 的一个给定\personDarboux{}--\personWeinstein{}邻域。若对所有
$t\in[0,1]$ 都有 $\phi_H^t(L)\subset V$，则称哈密顿同痕
$L_t=\phi_H^t(L)$ 为\termVEngulfedHamiltonianIsotopy{}。
\end{definition}

固定 $L$ 的\personDarboux{}邻域 $U$ 及其余切坐标卡
\[
  \Phi:U\longrightarrow V\subset T^*L,
\]
其中 $V$ 是 $T^*L$ 的零截面的一个邻域。若 $\phi$ 在 $C^1$ 意义下接近
$\mathrm{id}$，则可将 $\Phi(\phi(L))$ 写成 $\operatorname{Graph}df$，其中
$f:L\to\mathbb R$ 是光滑函数。我们假设 $\|df\|_{C^1}$ 充分小。

当 $\|df\|_{C^1}$ 变小时，\personFloer{}流线的模空间会出现\termThickThinDichotomy{}。这一二分在
(Oh96b) 中得到了详细分析，后来事实表明，它是拉格朗日\personFloer{}同调在辛拓扑中
应用的关键成分。出现这种二分，主要是因为薄流线与厚（或非薄）流线之间存在一个
“能隙”。

为精确定义这个间隙，对每个不依赖于时间的相容殆复结构
$J_0\in\mathcal J_\omega$，使用式 (8.4.32) 中引入的几何不变量
$A(\omega,L;J_0)$。回忆
\[
 A(\omega,J_0:L)=\min\{A_S(\omega,J_0),A_D(\omega,J_0;L)\}>0,
\]
其中
\[
 A_S(\omega,J_0):=\inf\left\{\int v^*\omega>0\ \middle|\
 v:S^2\to M,\ \overline\partial_{J_0}v=0\right\},
\]
而对给定的拉格朗日子流形 $L\subset(M,\omega)$，
\[
 A_D(\omega,J_0;L):=\inf\left\{\int w^*\omega>0\ \middle|\
 w:(D^2,\partial D^2)\to(M,L),\ \overline\partial_{J_0}w=0\right\}.
\]

可以把这个不变量推广到相容殆复结构的任意紧族
$K\subset\mathcal J_\omega$。从 $C^1$ 拓扑中的连续 $k$-参数族
$K:[0,1]^k\to\mathcal J_\omega$ 出发，定义常数
\[
 A_S(\omega;K)=\inf_{\kappa\in[0,1]^k}\{A_S(\omega,J(\kappa))\}.
\]
这个常数也是正的，并具有如下下半连续性。
\footnote{译者注：原书这里印作族内的上确界，但随后的下半连续命题及其反证都
使用族内能量下界，因而应为下确界；译文据此改正。}

\begin{proposition}
\label{prop:family-sphere-energy-lsc}
常数 $A_S(\omega;K)$ 关于 $K$ 下半连续。换言之，给定任意 $K$ 以及
$0<\epsilon<A_S(\omega;K)$，存在 $\delta=\delta(K,\epsilon)>0$，使得对任意
满足 $\|K'-K\|_{C^1}\leq\delta$ 的 $K'$，都有
\[
 A_S(\omega;K')\geq A_S(\omega;K)-\epsilon.
\]
\end{proposition}

\begin{proof}
用反证法。反设存在 $\epsilon>0$ 及序列 $K_i\to K$，使得对所有 $i$ 都有
\begin{equation}
 A_S(\omega;K_i)<A_S(\omega;K)-\epsilon.
\label{eq:17-family-energy-contradiction}
\end{equation}
\footnote{译者注：命题定义并讨论的是 $A_S(\omega;K)$，但原书证明在本式及
随后两处均省去下标 $S$，写作未在此定义的 $A(\omega;K_i)$；证明只使用全纯球面，
故译文补回下标 $S$。同段又把 $k$ 参数族的参数域印作 $[0,1]^n$，译文改为
$[0,1]^k$。}
对每个 $i$，由 $A_S(\omega;K_i)$ 的定义和~\eqref{eq:17-family-energy-contradiction}，
存在 $t_i\in[0,1]^k$ 及非常值 $J_{t_i}$-全纯球面 $v_i$，满足
\[
 A_S(\omega;K_i)\leq\omega([v_i])<A_S(\omega;K)-\frac{\epsilon}{3}.
\]
由\personGromov{}紧性，$v_i$ 收敛到一个 $J_{t_\infty}$-尖点曲线
\[
 v_\infty=\sum_a v_{\infty,a},
\]
其中 $J_{t_\infty}\in K$。若 $v_\infty$ 为常值，则当 $i$ 充分大时 $v_i$ 与零同调，
故 $\omega([v_i])=0$；这又蕴含 $v_i$ 是常值映射，与既定假设矛盾。若某个
$v_{\infty,a}$ 是非常值 $J_{t_\infty}$-全纯球面，则
\[
 \omega([v_{\infty,a}])\leq A_S(\omega;K)-\frac{\epsilon}{3},
\]
这与 $A_S(\omega;K)$ 的定义矛盾。证明完毕。
\end{proof}

类似论证证明不变量 $A(\omega,K;L)$ 的半连续性。

现在设 $U$ 是 $(M,\omega)$ 中 $L$ 的\personDarboux{}邻域，$H$ 是满足
\begin{equation}
 \phi_H^1(L)\subset U
\label{eq:17-hamiltonian-image-in-U}
\end{equation}
的哈密顿量，并令 $J=\{J_t\}_{t\in[0,1]}$ 是依赖于 $t$ 的殆复结构族。我们研究
记作 $\mathcal M(J;L,\phi_H^1(L))$ 的如下方程解的模空间：
\[
 \frac{\partial u}{\partial\tau}+J_t\frac{\partial u}{\partial t}=0,
 \quad u(\tau,0)\in L,\quad u(\tau,1)\in\phi_H^1(L),\quad E_J(u)<\infty.
\]
以 $d_{C^1}(L_0,L_1)$ 表示 $L_0$ 与 $L_1$ 的 $C^1$ 距离。

\begin{proposition}
\label{prop:nearby-thick-thin-gap}
给定任意 $0<\alpha<A(\omega,L;J_0)$ 及任意不依赖于时间的 $J_0$，存在常数
$\epsilon>0$，使得只要
$d_{C^1}(\phi_H^1(L),L)<\epsilon$ 且
$\max_{t\in[0,1]}|J_t-J_0|_{C^0}<\epsilon$，便有：
\begin{enumerate}[label=(\arabic*)]
\item 对所有满足 $\operatorname{Im}u\subset U$ 的
$u\in\mathcal M(J;L,\phi_H^1(L))$，都有
\begin{equation}
 \int u^*\omega<A(\omega,L;J_0)-\alpha,
\label{eq:17-thin-area-upper}
\end{equation}
反之亦然；
\item 其他所有 $u\in\mathcal M(J;L,\phi_H^1(L))$ 的像均不包含于 $U$，并且满足
\begin{equation}
 \int u^*\omega>A(\omega,L;J_0)-\delta,
\label{eq:17-thick-area-lower}
\end{equation}
其中 $\delta$ 依赖于 $\epsilon>0$（但不依赖于 $\alpha$），且充分小。
\end{enumerate}
\end{proposition}

\begin{proof}
先给出一个简单引理。

\begin{lemma}
\label{lem:strip-area-graph-df}
设 $L_0$ 是 $T^*L$ 的零截面，$L_1=\operatorname{Graph}df$，其中 $f$ 是 $L$ 上
的函数。假设 $L_0$ 与 $L_1$ 横截相交，且
$x,y\in L_0\cap L_1=\operatorname{Crit}(f)$。设
$u:[0,1]\times[0,1]\to T^*L$ 是 $C^1$ 映射，满足
\[
 u(s,0)\subset L_0,\quad u(s,1)\subset L_1,\quad
 u(0,t)\equiv y,\quad u(1,t)\equiv x.
\]
则
\[
 \int u^*\omega=f(y)-f(x).
\]
\footnote{译者注：原书在本引理边界条件中印作 $u(s,1)\subset L$；结合
$L_1=\operatorname{Graph}df$ 及证明第一句，译文改为 $L_1$。}
\end{lemma}

\begin{proof}
记 $u(s,1)=(q(s,1),p(s,1))$。由假设 $u(s,1)\in\operatorname{Graph}df$，故
$p(s,1)=df(q(s,1))$。定义
\[
 u_f(s,t)=(q(s,1),t\,df(q(s,1))).
\]
把 $u$ 与 $u_f$ 拼接得到 $w:=u\#u_f$，其边界包含在 $L_0$ 中。因此
\[
 0=\int w^*(-d\theta)=\int w^*\omega_0
   =\int u^*\omega-\int u_f^*\omega.
\]
另一方面，
\[
 \int u_f^*\omega=-\int_{\partial u_f}\theta
 =-\int_{u_f|_{t=1}}\theta
 =-\int_{s\mapsto q(s,1)}df=f(y)-f(x).
\]
\footnote{译者注：原书证明首句记 $u(1,t)=(q(1,t),p(1,t))$，$u_f$ 的定义亦随之
印作 $u_f(s,t)=(q(1,t),s\,df(q(1,t)))$。按引理的边界条件，含于
$\operatorname{Graph}df$ 的是 $t=1$ 边而非 $s=1$ 边，故此两处系槽位误植；译文
改记 $u(s,1)=(q(s,1),p(s,1))$，并把 $u_f$ 补正为
$(q(s,1),t\,df(q(s,1)))$。原书末式道路 $s\mapsto q(s,1)$ 与补正后的 $u_f$
一致，译文照录。}
证明完毕。
\end{proof}

当 $\phi_H^1(L)$ 在 $C^1$ 意义下接近 $L$ 时，对
$u\in\mathcal M^J(L,\phi_H^1(L);U)$ 应用引理，可写
$\Phi(\phi_H^1(L))=\operatorname{Graph}df$。由于 $\Phi$ 是辛的，即
$\Phi^*\omega_0=\omega$，故
\[
 \int u^*\omega=\int(\Phi\circ u)^*\omega_0=f(y)-f(x).
\]
此外，只要 $d_{C^1}(L,\phi_H^1(L))$ 充分小（例如小于 $\epsilon$），就有与 $H$
无关的 $C>0$，使得
\[
 \frac1C|df|_{C^0}\leq d_{C^1}(L,\phi_H^1(L))\leq C|df|_{C^0}.
\]
显然
\[
 |f(y)-f(x)|\leq\max f-\min f\leq |df|_{C^0}\operatorname{diam}(L).
\]

给定 $0<\alpha<A(\omega,L;J_0)$。先证明 (1)。若
$d_{C^1}(L,\phi_H^1(L))<\epsilon$ 且 $\epsilon$ 充分小，则
\[
 |df|_{C^0}<C\epsilon\leq
 \frac{A(\omega,L;J_0)-\alpha}{\operatorname{diam}(L)}.
\]
于是
\[
 \begin{aligned}
 \int u^*\omega&=f(y)-f(x)\leq\max f-\min f
 \leq |df|_{C^0}\operatorname{diam}(L)\\
 &\leq\frac{A(\omega,L;J_0)-\alpha}{\operatorname{diam}(L)}
 \operatorname{diam}(L)=A(\omega,L;J_0)-\alpha.
 \end{aligned}
\]

现证其逆命题。反设存在 $0<\alpha<A(\omega,L;J_0)$，以及序列
$(J^i,H^i)$，使得当 $i\to\infty$ 时
$d_{C^1}(\phi_{H^i}^1(L),L)\to0$ 且
$\max_{t\in[0,1]}|J_t^i-J_0|_{C^1}\to0$，并存在
$u_i\in\mathcal M(\phi_{H^i}^1(L),L;J^i)$，对所有 $i$ 满足
\footnote{译者注：原书已引入哈密顿量序列 $H^i$，但此处模空间中漏掉上标 $i$；
译文按同段的 $\phi_{H^i}^1(L)\to L$ 补正。}
\footnote{译者注：命题假设只要求 $J_t$ 与 $J_0$ 在 $C^0$ 意义下接近，但原书
反证在此改用 $C^1$ 收敛；译文照录这一加强的假设，因为原文未补出由 $C^0$ 控制
得到所需 $C^1$ 收敛的论证。}
\begin{equation}
 \operatorname{Im}u_i\not\subset U,
 \qquad \int u_i^*\omega\leq A(\omega,L;J_0)-\alpha.
\label{eq:17-thin-converse-sequence}
\end{equation}
由 $A(\omega,L;J_0)$ 的定义，此序列既不能冒出 $J_0$-全纯球面，也不能冒出圆盘。
由于 $\phi_{H^i}^1(L)\to L$（按 $C^1$ 拓扑）且 $L$ 紧，取子列后可设存在
$x_i,y_i\in\phi_{H^i}^1(L)\cap L$，满足 $x_i\to x_\infty$、
$y_i\to y_\infty$，其中 $x_\infty,y_\infty\in L$，并且
\[
 u_i\in\mathcal M(\phi_{H^i}^1(L),L;x_i,y_i;J^i).
\]
再由~\eqref{eq:17-thin-converse-sequence}，必要时沿 $\tau$ 方向平移 $u_i$，可设
\[
 u_i(0,t_i)\to p_\infty\in M\setminus U.
\]
序列不冒泡，故可应用局部收敛；得到局部极限
$J_0$-全纯映射 $u_\infty:\mathbb R\times[0,1]\to M$，满足
\[
 u_\infty(0,t_\infty)=p_\infty\notin U,\qquad
 u_\infty(\tau,0),u_\infty(\tau,1)\subset L,
 \qquad \int u_\infty^*\omega\leq A(\omega,L;J_0)-\alpha.
\]
利用 $\mathbb R\times i[0,1]$ 与 $D^2\setminus\{-1,1\}$ 的共形等价，可将
$u_\infty$ 看作定义在 $D^2\setminus\{-1,1\}$ 上的非常值映射，且
\[
 \overline\partial_{J_0}u_\infty=0,\quad
 u_\infty(\partial D^2\setminus\{-1,1\})\subset L,\quad
 \int u_\infty^*\omega<A(\omega,L;J_0)-\alpha.
\]
由可去奇点定理，它延拓为非常值 $J_0$-全纯映射
$u_\infty:(D^2,\partial D^2)\to(M,L)$，满足
\[
 0<\int u_\infty^*\omega\leq A(\omega,L;J_0)-\alpha,
\]
与 $A(\omega,L;J_0)$ 的定义矛盾，故 (1) 得证。

为证 (2)，仍作反设：存在 $\delta>0$、满足
$\phi_{H^k}(L)\to L$ 的 $H^k$、满足 $J^k\to J_0$（按 $C^1$）的 $J^k$，以及
$u_k\in\mathcal M(J^k;\phi_{H^k}^1(L),L)$，其像不包含于 $U$，但
\begin{equation}
 \int u_k^*\omega<A(\omega,L;J_0)-\delta.
\label{eq:17-thick-contradiction}
\end{equation}
在\personDarboux{}坐标卡中把 $\phi_{H^k}^1(L)$ 写成
$\operatorname{Graph}(df_k)$，只需考虑在 $U$ 中
$H_k=f_k\circ\pi$ 且 $|df_k|_{C^0}\to0$ 的情形。细化 (1) 的论证，在一致面积界
\eqref{eq:17-thick-contradiction} 下，得到仍记作 $u_k$ 的子列局部收敛到
$u_\infty$。因假设 $u_k$ 的像不包含于 $U$，$u_\infty$ 的像经过
$M\setminus U$ 中一点。由可去奇点定理，$u_\infty$ 给出一个边界落在 $L$ 上的
非常值圆盘或一个球面。因此按 $A(\omega,L;J_0)$ 的定义，
\[
 \int u_\infty^*\omega\geq A(\omega,L;J_0),
\]
从而
\[
 \liminf_{k\to\infty}\int u_k^*\omega
 \geq\int u_\infty^*\omega\geq A(\omega,L;J_0).
\]
但当 $k$ 充分大时，这与~\eqref{eq:17-thick-contradiction} 矛盾。证明完毕。
\end{proof}

\section{\termLocalFloerHomology{}}
\label{sec:local-floer-homology}

本节依照 (Fl89b)、(Oh96b) 的思路，定义\personFloer{}上同调 $HF(L,L)$ 的一个局部版本，
它只取像包含在给定\personDarboux{}--\personWeinstein{}邻域 $U$ 中的\personFloer{}流线所作的贡献。把
这个\termLocalFloerHomology{}记作 $HF(L,L;U)$。我们将证明，这个局部贡献只依赖于
$(L,U)$，因而可以在零截面 $o_L\cong L$ 的一个邻域
$V\subset T^*L$ 中，对 $(o_L,V)$ 作计算。此外，根据定义，
\stfhIndexRangeStart{I120}{Jubufuluoertongdiao}{局部弗洛尔同调（local Floer homology）}
$HF(L,L;U)$ 总是良定义的，并不需要 $L\subset M$ 满足诸如 $\omega$ 正合或
$(L,M)$ 单调之类的无障碍假设。

作为直接应用，我们依照 (Oh96b) 的方案证明\personGromov{}关于正合拉格朗日嵌入的两个
著名定理。

\subsection{局部\personFloer{}同调的定义}
\label{subsec:definition-local-floer}

我们从\personFloer{}在 (Fl89b) 的哈密顿情形中引入的如下灵活概念开始。它是动力系统中
\stfhIndex{I062}{Kangli}{康利（Conley）}
\personConley{} (Co78) 引入的\emph{孤立块}概念的\personFloer{}理论对应物。这个定义可以在更一般的抽象
框架中表述，但我们将集中于当前几何情形。

设 $U\subset M$ 是 $L$ 的一个\personDarboux{}邻域。引入如下定义。

\stfhIndex{I114}{Gulibubianji}{孤立不变集（isolated invariant set）}
\begin{definition}
\label{def:isolated-invariant-set}
设 $L'\subset U$ 是任意紧拉格朗日子流形，$J\in\mathcal J_\omega$。考虑带一个
标记点的\personFloer{}流线集合 $\mathcal M_1(L,L';J)$ 及其评价映射
$\mathcal M_1(L,L';J)\to M$。定义子集
\[
 \mathcal M(L,L';J;U):=\{u\in\mathcal M(L,L';J)\mid
 \operatorname{Image}u\subset U\}
\]
及其评价像
\[
 S(L,L';J;U)=\operatorname{ev}(\mathcal M_1(L,L';J;U)).
\]
称 $S(L,L';J;U)$ 为 $U$ 中的一个不变集。若
$\overline{S(L,L';J;U)}\subset U$，其中闭包在 $M$ 中取，则称它（在\personCauchy{}--\personRiemann{}流下）孤立于 $U$。当其孤立时，称
$\overline{S(L,L';J;U)}$ 为 $U$ 中的极大不变集。
\end{definition}

\stfhIndex{I067}{Yantuo}{延拓（continuation）}
现在定义极大不变集之间的延拓。

\begin{definition}
\label{def:continuation-maximal-invariant-sets}
设 $\{L^s\}_{s\in[0,1]}$ 是紧拉格朗日子流形 $L^0$ 的一个同痕。考虑
\[
 (\mathbb J,\mathbb H)\in
 \mathcal P([0,1]^2,\mathcal J_\omega)
 \times C^\infty([0,1]^2\times M,\mathbb R)
\]
以及开子集 $\mathbb U\subset[0,1]\times M$。以
$S_0:=S(J^0,(L,L^0);U^0)$ 和
$S_1:=S(J^1,(L,L^1);U^1)$ 表示上述极大不变集。若满足：
\begin{enumerate}[label=(\arabic*)]
\item 对每个 $s\in[0,1]$ 及所有 $t\in[0,1]$，
\[
 L^s\subset U^s:=\{x\in M\mid(s,x)\in\mathbb U\};
\]
\item 对所有 $s\in[0,1]$，每个不变集
\[
 S^s:=S(J^s,(L,L^s);U^s)
\]
都孤立于 $U^s$，
\end{enumerate}
\footnote{译者注：原书此式写作 $\{x\in P\mid(x,s)\in\mathbb U\}$，但 $P$
未定义，且与 $\mathbb U\subset[0,1]\times M$ 的坐标次序相反；译文据后者改正。}
则称 $(\mathbb J,\mathbb H,\mathbb U)$ 是极大不变集
$S_0\subset U^0$ 与 $S_1\subset U^1$ 之间的一个延拓。
\end{definition}

下面考虑按如下意义接近 $L$ 的拉格朗日子流形族。

\stfhSubIndex{I076}{Zhenghe}{正合（exact）}{XiangduiyuLzhenghe}{相对于 $L$ 正合（exact relative to $L$）}
\begin{definition}
\label{def:exact-relative-to-L}
若存在包含 $L$ 的\personDarboux{}邻域 $U$，使得在
$U\cong V\subset T^*L$ 中 $L'$ 是正合的，则称拉格朗日子流形
$L'\subset(M,\omega)$ \termExactRelativeToL{}。
\end{definition}

有了这些定义，下面的结论容易证明。

\begin{lemma}
\label{lem:regular-perturbation-isolated-set}
设 $L\subset M$，且 $L'$ 相对于 $L$ 正合并与 $L$ 横截相交。假设
$S(L,L';J;U)$ 孤立于 $U$。则存在 $J$ 的一个 $C^\infty$ 扰动 $J'$，使得
$\mathcal M(J',L,L';U)$ 是\personFredholm{}正则的，且 $S(J',L,L';U)$ 仍孤立于
$U$。特别地，对任意满足 $\mu(x;U)-\mu(y;U)=1$ 的
$x,y\in L\cap L'$，
$\mathcal M(x,y;J';U)\subset\mathcal M(L,L';J';U)$ 的基数有限。
\end{lemma}

\begin{proof}
按定义，孤立性在 $C^\infty$ 扰动下稳定，故这样的扰动 $J'$ 立即存在。又因
$(L,L')$ 在 $U$ 中不会冒泡，紧性给出有限性。
\end{proof}

设 $(L,L';J;U)$ 和 $J'$ 如引理 17.2.4 所述。定义
\[
 n_U(x,y;J'):=\#\{\mathcal M((x,y);J';U)\text{ 中的孤立流线}\}\pmod2.
\]
若 $L'=\phi_H^1(L)$，则利用第 15.6 节建立的相容定向，可以定义整数
$n_U(x,y;J')$。

\begin{theorem}
\label{thm:local-floer-boundary-continuation}
假设 $(L,L';J;U)$ 如引理 17.2.4 所述。对任意使
$\mathcal M(L,L';J';U)$ \personFredholm{}正则的 $J$ 的小扰动 $J'$，同态
\[
 \partial_U:CF(L,L';J';U)\longrightarrow CF(L,L';J';U),
 \qquad
 \partial_Ux=\sum_{y\in L\cap L'}n_U(x,y;J')y
\]
满足 $\partial_U\circ\partial_U=0$。此外，只要延拓
$(S,\mathbb J,\mathbb H,\mathbb U)$ 在端点 $s=0,1$ \personFloer{}正则，相应商空间
\[
 HF^*(L,L';J';U)=\ker\partial_U/\operatorname{im}\partial_U
\]
在该延拓下同构。
\footnote{译者注：定理前半以一般 $L'$ 为对象，但原书求和指标印作
$y\in L\cap\phi_H^1(L)$；译文按定理设定改为 $y\in L\cap L'$。}
\end{theorem}

\begin{proof}
关于边界映射的断言来自孤立性假设，因为三元组 $(L,L';U)$ 不承载任何气泡。
对延拓不变性，我们取单位区间的一个分割
$0=t_0<t_1<\cdots<t_N=1$，并把与之相伴的一列“绝热链映射”复合起来，构造两端
之间的链映射。因此在余下证明中，可以假设两端之间的 $C^1$ 距离任意小。

考虑“非自治”\personFloer{}流线方程
\begin{equation}
 \left\{
 \begin{aligned}
 &\partial u/\partial\tau+J_t^{\rho(\tau)}\partial u/\partial t=0,\\
 &u(\tau,0)\in L^{\rho(\tau)},\qquad u(\tau,1)\in L,
 \end{aligned}
 \right.
\label{eq:17-nonautonomous-floer}
\end{equation}
其中 $\mathbb J=\{J^s\}$ 是两个\personFredholm{}正则结构 $J^0,J^1$ 之间的同伦，
$\mathbb L=\{L^s\}$ 是 $L^0$ 与 $L^1$ 之间的哈密顿同痕。以
$\mathcal M(L,L^\rho;J^\rho)$ 表示~\eqref{eq:17-nonautonomous-floer} 的解的模空间，
以 $\overline{\mathcal M}(L,L^\rho;J^\rho)$ 表示其闭包；后者由第 14.3 节所定义的
稳定断裂\personFloer{}流线组成。

定义
\begin{equation}
 d_{C^1}(L',\mathbb L):=\max_{s\in[0,1]}d_{C^1}(L',L^s),
 \qquad
 d_{C^0}(J_0,\mathbb J):=\max_{(s,t)\in[0,1]^2}|J_0-J_t^s|_{C^0}.
\label{eq:17-continuation-distances}
\end{equation}
令 $L^s=\phi_{H_i}^s(L')$。只需证明如下结论。

\begin{lemma}
\label{lem:continuation-evaluation-in-U}
考虑评价映射
$\overline{\operatorname{ev}}^\rho:
\overline{\mathcal M}(L,L^\rho;J^\rho)\to M$。存在
$\epsilon>0$，使得若
$d_{C^1}(L',\mathbb L)<\epsilon$ 且
$d_{C^0}(J_0,\mathbb J)<\epsilon$，则
\[
 \operatorname{Image}\overline{\operatorname{ev}}^\rho\subset U.
\]
\footnote{译者注：原书引理条件中第二个距离的首项印作 $J$，即
$d_{C^0}(J,\mathbb J)$，但式 (17.2.8) 定义的是
$d_{C^0}(J_0,\mathbb J)$；译文按该定义改正。}
\end{lemma}

\begin{proof}
我们知道
\[
 \mathcal M(L,L^\rho;J^\rho)
 =\bigcup_{x\in L\cap L^0,\ y'\in L\cap L^1}
 \overline{\mathcal M}(x,y';J^\rho).
\]
反设存在 $\mathbb J_i=\{J_i^s\}$ 和 $H_i$，其中
$L^s=\phi_{H_i^s}(L')$，且当 $i\to\infty$ 时
$d_{C^1}(L,\phi_{H_i}^1(L'))\to0$、
$d_{C^0}(J,\mathbb J_i)\to0$，并存在
$u_i\in\mathcal M(x,y;L,L^\rho;J^\rho)$ 及点 $(\tau_i,t_i)$，使得
\begin{equation}
 u_i(\tau_i,t_i)\in M\setminus U.
\label{eq:17-continuation-escapes-U}
\end{equation}
取子列后可设 $\rho(\tau_i)\to s_0$、$t_i\to t_0$。利用能量界与条件
\[
 d_{C^1}(L,\mathbb L)\to0,
 \qquad d_{C^0}(J,\mathbb J)\to0,
\]
并适当沿 $\tau$ 方向平移序列，可以抽取一个满足关于 $J_0$ 的方程
\eqref{eq:17-nonautonomous-floer} 的局部极限 $u_\infty$，使得对某些
$x_0,y_0\in L\cap L'$，
\[
 u_\infty(-\infty)\to x_0,\quad u_\infty(\infty)\to y_0,\quad
 u_\infty(0,t_0)\in M\setminus U.
\]
这与 $S(L,L';J;U)$ 孤立于 $U$ 的假设矛盾。
\end{proof}

对引理中的同伦 $(\mathbb J,\{L^s\})$，可定义
\[
 n_U(p,q';J^\rho)=\#\mathcal M(p,q';J^\rho;U).
\]
于是同态
\[
 h_{(J^\rho,L^\rho;U)}:CF(L,L^0;U)\longrightarrow CF(L,L^1;U)
\]
良定义，且对任意\personFloer{}正则的 $\mathbb J$ 满足
\[
 h_{(J^\rho,L^\rho;U)}\partial_{(L,L^0;J^0;U)}
 +\partial_{(L,L^1;J^1;U)}h_{(J^\rho,L^\rho;U)}=0.
\]
这就建立了 $h_{(J^\rho,L^\rho;U)}$ 的链性质。用与第 12.5.2 节相同的论证可证明
它在同调上诱导同构，从而完成定理 17.2.5 的证明。
\footnote{译者注：原书此处印作“定理 17.2.7”，但本段属于定理 17.2.5 的证明，
且真正的定理 17.2.7 尚在后文出现；译文据此改正。}
\end{proof}

这就完成了任意相对于 $L$ 正合的 $L'$ 的局部\personFloer{}上同调
$HF(L,L';U)$ 的构造。

\subsection{吞没哈密顿同痕}
\label{subsec:engulfed-hamiltonian-isotopy}

本节余下部分考虑给定紧拉格朗日子流形 $L\subset(M,\omega)$ 的\termLocalFloerHomology{}得以定义的自然情形。设 $V\subset\overline V\subset U$ 是 $L$ 的一对给定\personDarboux{}邻域。

先在 $U$ 中以孤立延拓，把 $(L,\phi_H^1(L);J;U)$ 的极大不变集形变到
$(L,\phi_H^1(L);J_0;U)$ 的极大不变集；这里 $J_0$ 不依赖于时间，并且在
$U\subset T^*L$ 上限制为与 $L$ 上\personRiemann{}度量 $g$ 相伴的典范殆复结构 $J_g$。
（$J_g$ 的定义见定义 12.2.1。）引理 17.2.6 保证这种形变是可能的。

再用 (Fl89a) 中的论证，即可完成如下计算。

\begin{theorem}
\label{thm:local-floer-is-morse}
设 $L\subset M$ 如上，$U$ 是 $L$ 的一个\personDarboux{}邻域。若对某个不依赖于时间的
$J_0$，有
$d_{C^1}(L,\phi_H^1(L))\leq\epsilon_3$、
$d_{C^0}(J_0,J)<\epsilon_3$，且 $J$ 是 $(L,\phi_H^1(L))$-正则的，则
\[
 HF((L,\phi_H^1(L));J;U)\cong H_*(L;\mathbb Z_2).
\]
\end{theorem}

\begin{proof}
选择一个可以显式计算 $HF(J,(L,\phi_H^1(L));U)$ 的特殊正则对
$(J,\phi_H)$。为此，利用 $U$ 辛同胚于 $T^*L$ 零截面的一个邻域，把一个
$C^2$-小\personMorse{}函数 $f$ 的图像 $\operatorname{Graph}df\subset T^*L$ 植入
$U$。

取零截面的邻域 $V$，满足 $V\subset\overline V\subset U$。考虑与 $g$ 相伴的
典范殆复结构 $J_g$，以及自治哈密顿量 $H=f\circ\pi$ 诱导的哈密顿同痕
$\phi^{f\circ\pi}:T^*L\to T^*L$，其中 $f$ 是 $L$ 上的\personMorse{}函数，
$\pi:T^*L\to L$ 是典范投影。一般而言，$\phi^{f\circ\pi}$ 未必把 $U$ 映到
$U$；但若 $f$ 在 $C^1$ 意义下充分小，则对所有 $t\in[0,1]$，
\[
 \phi_{f\circ\pi}^t(\overline V)\subset U.
\]
对每个光滑映射 $\chi:\mathbb R\to L$，定义映射
$u_\chi:\mathbb R\times[0,1]\to P$：
\[
 u_\chi(\tau,t)=\phi_{f\circ\pi}^t(\chi(\tau)),
 \qquad (\tau,t)\in\mathbb R\times[0,1].
\]
仍在 $f$ 为 $C^1$-小的条件下，对所有 $t\in[0,1]$ 有
$\phi_{f\circ\pi}^t(L)\subset V$，从而
\[
 \operatorname{Im}u_\chi\subset V\subset U.
\]
在 $T^*L$ 上定义殆复结构族 $J^f=\{J_t^f\}$：
\[
 J_t^f=(\phi_{f\circ\pi}^t)_*J_g(\phi_{f\circ\pi}^t)_*^{-1}.
\]
让 $f$ 在 $C^1$ 意义下变小，就能使它任意接近不依赖于时间的 $J_g$。
最后对余切丛 $T^*N$ 中的 $L=o_N$ 应用推论 12.5.7，即得结论。
\end{proof}

值得指出，若 $L$ 是正合 $(M,\omega)$ 中的紧正合拉格朗日子流形，则
$A(\omega,J;L)=\infty$。即使非紧正合辛流形 $M$ 在无穷远处驯服或有界，这一
结论仍成立。于是立即得到：

\begin{theorem}
\label{thm:exact-lagrangian-floer-homology}
设 $(M,\omega)$ 是驯服的正合辛流形，$L\subset(M,\omega)$ 是紧正合拉格朗日
子流形。则
\[
 HF_*(L,L)\cong H_*(L).
\]
\end{theorem}

\begin{proof}
固定 $L$ 的\personDarboux{}--\personWeinstein{}邻域 $U$。取\personMorse{}函数 $f:L\to\mathbb R$，
并令 $\epsilon>0$ 充分小。由于 $A(\omega,J;L)=\infty$，命题 17.1.3 中只有情形
(1) 发生：当 $\epsilon>0$ 充分小且 $|J-J_0|_{C^0}\leq\epsilon$ 时，
$\mathcal M(L,\phi_{\epsilon f}(L);J)$ 中所有流线都包含在 $U$ 中。因此
\[
 HF_*(L,\phi_{\epsilon f}(L))
 \cong HF_*(L,\phi_{\epsilon f}(L);U)
 \cong H_*(L).
\]
\end{proof}

这个定理立即给出\personGromov{} (Gr85) 的以下两个定理；正如 (Oh96b) 所展示的，它们
揭示了环境辛流形的两种不同特征。

\begin{theorem}[\personGromov{}]
\label{thm:no-exact-lagrangian-in-Cn}
$\mathbb C^n$ 中不存在正合拉格朗日嵌入。
\end{theorem}

\begin{proof}
若存在这样的嵌入 $L\subset\mathbb C^n$，则它的\personFloer{}同调（例如取
$\mathbb Z_2$ 系数）同构于 $H_*(L;\mathbb Z_2)$。由于 $L$ 紧，至少有
$H_0(L;\mathbb Z_2)\neq0$，故对任意哈密顿同痕 $\phi_H$，
$HF_*(L,\phi_H^1(L);J;\mathbb Z_2)\neq0$。特别地，对任意（不一定横截的）
哈密顿微分同胚 $\phi_H^1$，都有
\[
 L\cap\phi_H^1(L)\neq\varnothing.
\]
另一方面，显然可以用一个本身为哈密顿同痕的平移把 $L$ 从自身移开，矛盾。
\end{proof}

这次考虑任意紧流形 $N$ 的余切丛 $T^*N$。

\begin{theorem}[\personGromov{}]
\label{thm:exact-lagrangian-meets-zero-section}
设 $L\subset T^*N$ 是任意紧正合拉格朗日嵌入。则
$L\cap o_N\neq\varnothing$。
\end{theorem}

\begin{proof}
仍用反证法。设 $L\cap o_N=\varnothing$。考虑径向数乘映射
\[
 R_\lambda:T^*N\to T^*N,\qquad(q,p)\longmapsto(q,\lambda p),\qquad\lambda>0.
\]
由 $L\cap o_N=\varnothing$，当 $\lambda>0$ 充分大时，
\begin{equation}
 R_\lambda(L)\cap L=\varnothing.
\label{eq:17-radial-displacement}
\end{equation}
$R_\lambda$ 并非辛映射，而只是共形辛映射：
$R_\lambda^*\omega_0=\lambda\omega_0$。然而，把 $R_\lambda$ 的共形辛性质与
$L$ 的正合性结合起来可知，同痕
\[
 \lambda\longmapsto R_\lambda(L)
\]
是正合拉格朗日同痕。由定理 3.6.7，可找出哈密顿量 $H$，使得
$\phi_H^\lambda(L)=R_\lambda(L)$；因此对任意 $\lambda$ 都有
$R_\lambda(L)\cap L\neq\varnothing$，再次矛盾。
\end{proof}

\stfhIndex{I094}{Shengu-Saideer-Shimisi}{深谷--塞德尔--史密斯（Fukaya--Seidel--Smith）}
\stfhIndex{I135}{Nadele}{纳德勒（Nadler）}
近来\personNadler{} (Na09) 以及\personFukaya{}、\personSeidel{}、\personIvanSmith{} (FSS08) 借助\termFukayaCategory{}的同调
工具，大大加强了定理 17.2.10。

\begin{remark}
\label{rem:c0-localization-engulfed-H}
在 (Oh11b) 对拓扑哈密顿路径的谱不变量的研究中，理解时间一映射在 $C^0$ 意义下
充分小的 $V$-吞没哈密顿量 $H$ 的\personFloer{}复形局部化非常重要。对这样的哈密顿量，
第 17.1 节通过考察面积和作用量来建立\termThickThinDichotomy{}的方法并不合适，因为当
$\phi_H^1$ 仅在 $C^0$ 意义下小时，我们无法控制对子
$(\phi_H^1(L),L)$ 的面积或作用量。事实表明，获得\personFloer{}模空间厚--薄分解的正确
方法，是利用殆复结构 $J_g$ 的\emph{最大值原理}。当 $H$ 在 $C^2$ 意义下小时，这一分解
与第 17.1 节给出的分解一致。详见 (Oh11b)。
\end{remark}
\stfhIndexRangeEnd{I120}{Jubufuluoertongdiao}{局部弗洛尔同调（local Floer homology）}

\section{谱序列的构造}
\label{sec:construction-spectral-sequence}

\stfhIndexRangeStart{I179}{Puxulie}{谱序列（spectral sequence）}
本节介绍一个重要工具——谱序列。它对计算\personFloer{}同调或把\personFloer{}同调用于实际问题
都很有用。这个特定谱序列把拉格朗日子流形 $L$ 的内蕴拓扑与嵌入
$L\subset(M,\omega)$ 的外在辛拓扑联系起来；作者在 (Oh96b) 中用小
$\epsilon>0$ 时的\personFloer{}复形
\[
 CF(\phi_{-\epsilon f}^1(L),L)
\]
（取 $\mathbb Z_2$ 系数）表述了它。\personBuhovsky{} (Bu10)、\personBiran{}--\personOctavianCornea{} (BCo09)
和\personDamian{} (D09) 仍以小\personMorse{}函数的\personMorse{}同调为基础，加强了这一构造。不过，
这个谱序列的最佳设定，是 (FOOO09) 通过直接考察干净对子 $(L,L)$ 而给出的\personBott{}--\personMorse{}设定。后一构造所需数学超出本书框架，故我们只讨论 (Oh96b) 的设定；
但至少会说明一个重要成分，即谱序列乘积结构的相容性。这里借用 (FOOO09) 定理 D
的叙述，并作适合单调情形的简化。

由于\personFloer{}同调中要用到乘积结构，本节采用\personFloer{}同调的上同调版本；这意味着把
$-\epsilon f$ 换成 $\epsilon f$。

基于\termThickThinDichotomy{}和局部\personFloer{}上同调的分析，固定\personMorse{}函数
$f:L\to\mathbb R$ 和充分小的 $\epsilon>0$，使得命题 17.1.3 与定理 17.2.7 的
二分对 $(L,\phi_{\epsilon f}^1(L))$ 成立。以 $\Sigma_L\geq2$ 表示 $L$ 的最小
正马斯洛夫数。

置 $N_L=[(n+1)/\Sigma_L]$，并如前定义形式参数
\[
 q^\beta=T^{\omega(\beta)}e^{\mu_L(\beta)},
 \qquad
 \deg(q^\beta)=\mu_L(\beta),\quad E(q^\beta)=\omega(\beta),
\]
再考虑\personNovikov{}环 $\Lambda_{(\omega,L)}$ 及其子环
$\Lambda_{0,(\omega,L)}$。

考虑\personFloer{}复形 $CF(L,\phi_{\epsilon f}^1(L))$，并按第 13.9.3 节的方式赋予它
$\Lambda_{(L,\omega)}$-模结构。回忆，第 13.9.3 节把与完整道路空间
$\Omega(L_0,L_1)$ 相伴的
$CF^*(L_1,L_0)\cong CF_*(L_0,L_1)^*$ 看成 $\Lambda_{(L_0,L_1)}$-模。

在当前 $(L,\phi_{\epsilon f}^1(L))$ 情形中，把 $q^\beta[p]$ 的滤过层级与分次定义为
\begin{equation}
 \ell(q^\beta[p])=\omega(\beta)+f(p),
\label{eq:17-filtration-q-beta-p}
\end{equation}
\begin{equation}
 \deg(q^\beta[p])=n-\mu_f(p)-\mu_L(\beta).
\label{eq:17-grading-q-beta-p}
\end{equation}
这样设定，是为了让 $\overline\delta_{(0)}$ 与 $-f$ 的\personMorse{}边界算子（或 $f$ 的\personMorse{}上边界算子）一致。换言之，在 $\Lambda_{(\omega,L)}$ 中延拓关系
\begin{equation}
 [\beta]\cdot[p]=T^{-\omega(\beta)}e^{-\mu_L(\beta)}\cdot[p],
\label{eq:17-module-relation-beta-p}
\end{equation}
由此给 $CF(L,\phi_{\epsilon f}^1(L))$ 赋予 $\Lambda_{(\omega,L)}$-模结构。按定义，
\[
 CF(L,\phi_{\epsilon f}^1(L))=\Lambda_{\omega,L}\{\operatorname{Crit}f\}.
\]
于是上边界映射
\[
 \delta:CF(L,\phi_{\epsilon f}^1(L))\longrightarrow
 CF(L,\phi_{\epsilon f}^1(L))
\]
保持这个 $\Lambda_{(\omega,L)}$-模结构。

由\termThickThinDichotomy{}，边界算子 $\delta:C^*\to C^*$ 具有形式
\[
 \delta=\delta_{(0)}+\delta',
\]
其中 $\delta_{(0)}$ 来自与局部\personFloer{}复形相伴的上边界算子，$\delta'$ 是厚流线的
贡献。

为进一步分析 $\delta'$，设
$u\in\mathcal M(L,\phi_{\epsilon f}^1(L);x,y)$ 是任意厚流线，
$w:([0,1]^2,\partial[0,1]^2)\to(M,L)$ 是引理 17.1.4 证明中使用的相伴映射。
若 $\epsilon>0$ 取得充分小，则
\[
 \int w^*\omega\geq\Sigma_L-\frac14\Sigma_L
 =\frac34\Sigma_L>0,
 \qquad
 \mu_L(w)=1-\mu_f(x)+\mu_f(y)\leq n+1.
\]
\footnote{译者注：原书在已把 $\Sigma_L$ 定义为最小马斯洛夫数后，直接在面积
不等式中用它与 $\int w^*\omega$ 比较。除非另有未声明的单调常数归一化，否则两边
量纲不合；因所缺比例常数未在此固定，译文照录而不作不可唯一确定的修补。}
由 $L$ 的单调性与 $\int w^*\omega>0$，有 $\mu_L(w)>0$。因而
\[
 \mu_f(x)-\mu_f(y)=1-\mu_L(w),
\]
且
\[
 \mu_L(w)=l\Sigma_L,
 \qquad 1\leq l\leq\left[\frac{n+1}{\Sigma_L}\right].
\]
所以，每条厚流线都把\personMorse{}分次改变 $-l\Sigma_L+1$，其中
$1\leq l\leq[(n+1)/\Sigma_L]$；于是 $\delta$ 分解为
\begin{equation}
 \delta'=\overline\delta_{(1)}\otimes q
 +\cdots+\overline\delta_{(N_L)}\otimes q^{N_L},
\label{eq:17-delta-prime-decomposition}
\end{equation}
这里 $\overline\delta_{(l)}:C^{(*)}(f)\to C^{(*-l\Sigma_L+1)}(f)$ 是由连接
指标为 $*$ 与 $*-l\Sigma_L+1$ 的 $f$ 的临界点的流线所诱导的线性映射。回忆
$\delta_{(0)}$ 把 $C^*$ 映到 $C^{*+1}$。
\footnote{译者注：原书把目标次数印作 $*-1+l\Sigma_L$，与上一句给出的分次变化
$-l\Sigma_L+1$ 以及 $l=0$ 时 $\delta_{(0)}:C^*\to C^{*+1}$ 均矛盾；译文据这
两处改为 $*-l\Sigma_L+1$。}
\footnote{译者注：原书先把 $\delta'$ 定义为厚流线的贡献，但公式 (17.3.14) 的
右端又从局部项 $\overline\delta_{(0)}$ 开始；随后按 $q$ 次数展开的论证唯一确定
厚项应从 $l=1$ 开始，译文据此删去该局部项。}

\begin{theorem}
\label{thm:lagrangian-spectral-sequence}
存在具有如下性质的谱序列：
\begin{enumerate}[label=(\arabic*)]
\item
\[
 E_2^{p,q}=\bigoplus_k H^k(L;\mathbb Q)\otimes
 \left(T^{q\lambda}\Lambda_{(\omega,L)}
 /T^{(q+1)\lambda}\Lambda_{(\omega,L)}\right)^{(p-k)},
\]
其中 $\lambda>0$；
\item \personFloer{}上同调 $HF(L,\phi_{\epsilon f}^1(L);\Lambda_{(\omega,L)})$ 上存在滤过
$F^*HF(L,\phi_{\epsilon f}^1(L);\Lambda_{(\omega,L)})$，使得
\[
 E_\infty^{p,q}\cong
 \frac{F^qHF^p(L,\phi_{\epsilon f}^1(L);\Lambda_{(\omega,L)})}
 {F^{q+1}HF^p(L,\phi_{\epsilon f}^1(L);\Lambda_{(\omega,L)})}.
\]
\end{enumerate}
\end{theorem}

（更一般的陈述见 (FOOO09) 定理 D，尤其是由 \emph{bulk 形变}或由 $M$ 的环境上同调类
形变的版本。）本小节余下部分证明这个定理。

因为 $\Sigma_L\geq2$，$\overline\delta_{(l)}$ 的次数不大于 $-1$。比较方程
\[
 0=\delta\circ\delta=(\delta_{(0)}+\delta')^2
\]
各项的次数，立即得到
\[
 \delta_{(0)}\circ\delta_{(0)}=0,
 \qquad
 \delta_{(0)}\circ\delta'+\delta'\circ\delta_{(0)}
 +\delta'\circ\delta'=0.
\]
第一式使我们可以定义
$H(C_*,\delta_{(0)})\cong H^*(L;f)$。利用
\eqref{eq:17-delta-prime-decomposition} 展开第二式：
\[
 \begin{aligned}
 0&=\delta_{(0)}\circ\delta'+\delta'\circ\delta_{(0)}
 +\delta'\circ\delta'\\
 &=\bigl(\delta_{(0)}\overline\delta_{(1)}
 +\overline\delta_{(1)}\delta_{(0)}\bigr)q\\
 &\quad+\bigl(\overline\delta_{(1)}\overline\delta_{(1)}
 +\delta_{(0)}\overline\delta_{(2)}
 +\overline\delta_{(2)}\delta_{(0)}\bigr)q^2+\cdots.
 \end{aligned}
\]
收集 $q$-次数为 1 和 2 的项，得到
\[
 \delta_{(0)}\overline\delta_{(1)}+
 \overline\delta_{(1)}\delta_{(0)}=0,
\]
\[
 \overline\delta_{(1)}\overline\delta_{(1)}
 +\delta_{(0)}\overline\delta_{(2)}
 +\overline\delta_{(2)}\delta_{(0)}=0.
\]
第一式说明 $\overline\delta_{(1)}$ 下降到
$H(C_*,\delta_{(0)})=H^*(L;\epsilon f)$ 上，第二式说明它在那里定义一个微分。

$C(L,\phi_{\epsilon f}^1(L))$ 上的作用量滤过诱导出所定义的谱序列。首先有同构
\[
 \operatorname{gr}_*CF(L,\phi_{\epsilon f}^1(L))
 \cong C(L;\epsilon f)\otimes
 \operatorname{gr}_*(F^*\Lambda_{(\omega,L)})
\]
作为 $\operatorname{gr}_*(F^*\Lambda_{(\omega,L)})$-模。（关于
$\Lambda_{(\omega,L)}$ 上的滤过，见定义 16.2.3 及其后。）于是得到谱序列
$E_r^{p,q}$，其中
\begin{equation}
 E_2^{p,q}=\bigoplus_{k=0}^n H^k(L;R)\otimes
 \operatorname{gr}_q\bigl(F\Lambda_{(\omega,L)}^{(p-k)}\bigr).
\label{eq:17-spectral-sequence-E2}
\end{equation}
这里 $\Lambda_{(\omega,L)}^{(p-k)}$ 表示次数 $p-k$ 的部分。还可证明
$HF^p(L,\phi_{\epsilon f}^1(L);\Lambda_{(\omega,L)})$ 上存在滤过，使得
\[
 E_\infty^{p,q}\cong
 F^qHF^p(L,\phi_{\epsilon f}^1(L);\Lambda_{(\omega,L)})
 /F^{q+1}HF^p(L,\phi_{\epsilon f}^1(L);\Lambda_{(\omega,L)}).
\]
上述谱序列的收敛 $E_2^{p,q}\Longrightarrow E_\infty^{p,q}$ 的证明对任意交换环
$R$ 都成立。

现在考虑带滤过 $F^*D$ 的分次 $\Lambda_{(\omega,L)}$-模 $D$。该滤过与
$\Lambda_{(\omega,L)}$ 上的滤过相容；对
$D=CF(L,\phi_{\epsilon f}^1(L))$，即
\[
 F^l\Lambda_{(\omega,L)}\cdot F^kD\subseteq F^{k+l}D.
\]
令 $\Lambda_{(\omega,L)}^{(0)}$ 为 $\Lambda_{(\omega,L)}$ 的零次部分。则
\[
 \Lambda_{(\omega,L)}^{(0)}
 \cong R[[T^{A(\omega,L)}]][T^{-A(\omega,L)}].
\]
这说明，若 $R$ 是域，则 $\Lambda_{(\omega,L)}^{(0)}$ 也是域，因此 $D$ 的 $p$ 次
部分 $D^p$ 是 $\Lambda_{(\omega,L)}^{(0)}$ 上的向量空间。依照滤过复形的标准过程
(McC85)，用这个滤过定义谱序列。

注意，分别乘以 $T^{A(\omega,L)}$ 与 $e$ 给出同构
\[
 T^{A(\omega,L)}\cdot:F^qHF^p(L,\phi_{\epsilon f}^1(L);\Lambda_{(\omega,L)})
 \longrightarrow F^{q+1}HF^p(L,\phi_{\epsilon f}^1(L);\Lambda_{(\omega,L)}),
\]
\[
 e\cdot:F^qHF^p(L,\phi_{\epsilon f}^1(L);\Lambda_{(\omega,L)})
 \longrightarrow F^qHF^{p+2}(L,\phi_{\epsilon f}^1(L);\Lambda_{(\omega,L)}).
\]
记
\[
 \Lambda^{(0)}(\lambda)=\Lambda_{(\omega,L)}^{(0)}
 /F^\lambda\Lambda_{(\omega,L)}^{(0)}.
\]
以
$F^n\Lambda_{(\omega,L)}^{(0)}=F^{nA(\omega,L)}\Lambda_{(\omega,L)}^{(0)}$
定义 $\Lambda_{(\omega,L)}^{(0)}$ 的滤过。其相伴分次模为
\[
 \operatorname{gr}_*\bigl(F\Lambda_{(\omega,L)}^{(0)}\bigr)
 =\bigoplus_{n\in\mathbb Z_{\geq0}}
 \operatorname{gr}_n\bigl(F\Lambda_{(\omega,L)}^{(0)}\bigr),
\]
其中每个 $\operatorname{gr}_n(F\Lambda_{(\omega,L)}^{(0)})$ 自然同构于
$\Lambda^{(0)}(A(\omega,L))$。又有
\[
 \operatorname{gr}_*(F\Lambda_{(\omega,L)})
 =\operatorname{gr}_*(F\Lambda_{(\omega,L)}^{(0)})[e,e^{-1}]
 =\bigoplus_{n\in\mathbb Z_{\geq0}}
 \operatorname{gr}_n(F\Lambda_{(\omega,L)}^{(0)})[e,e^{-1}].
\]

以 $\widehat C$ 表示 $\operatorname{Crit}f$ 上的
$\Lambda_{(\omega,L)}^{(0)}$-模，并定义
\[
 \begin{aligned}
 Z_r^{p,q}(\widehat C)
 &=\{x\in F^q\widehat C^p\mid
 \delta x\in F^{q+r-1}\widehat C^{p+1}\}+F^{q+1}\widehat C^p,\\
 B_r^{p,q}(\widehat C)
 &=\bigl(\delta(F^{q-r+2}\widehat C^{p-1})\cap F^q\widehat C^p\bigr)
 +F^{q+1}\widehat C^p,\\
 E_r^{p,q}(\widehat C)&=Z_r^{p,q}(\widehat C)/B_r^{p,q}(\widehat C).
 \end{aligned}
\]
回忆 $E_r^{p,q}$ 自然是 $\Lambda_{(\omega,L)}^{(0)}$-模。乘以
$e^{\pm1}\in\Lambda_{(\omega,L)}$ 定义映射
\[
 e^{\pm1}:E_r^{p,q}(\widehat C)\to E_r^{p\pm2,q}(\widehat C),
\]
从而使
$E_r(\widehat C):=\bigoplus_{p,q}E_r^{p,q}(\widehat C)$ 成为
$\operatorname{gr}_*(F\Lambda_{(\omega,L)})$-模，并且
\[
 \bigoplus_{p,q}\delta_r^{p,q}:
 \bigoplus_{p,q}E_r^{p,q}(\widehat C)
 \longrightarrow
 \bigoplus_{p,q}E_r^{p+1,q+r-1}(\widehat C)
\]
是 $\operatorname{gr}(F^*\Lambda_{(\omega,L)})$-模同态。

\begin{lemma}
\label{lem:spectral-sequence-differential}
设 $E_r^{p,q}(\widehat C)$ 是上面定义的 $\Lambda_{(\omega,L)}^{(0)}$-模。存在
$\Lambda_{(\omega,L)}^{(0)}$-模同态
\[
 \delta_r^{p,q}:E_r^{p,q}(\widehat C)\to
 E_r^{p+1,q+r-1}(\widehat C),
\]
满足
\begin{enumerate}[label=(\roman*)]
\item $\delta_r^{p+1,q+r-1}\circ\delta_r^{p,q}=0$；
\item $\ker(\delta_r^{p,q})/\operatorname{im}(\delta_r^{p-1,q-r+1})
 \cong E_{r+1}^{p,q}(\widehat C)$；
\item $e^{\pm1}\circ\delta_r^{p,q}
 =\delta_r^{p\pm2,q}\circ e^{\pm1}$。
\end{enumerate}
\end{lemma}

\begin{proof}
定义 $\delta_r^{p,q}[x]=[\delta x]\in E_r^{p+1,q+r-1}(\widehat C)$，结论立即成立。
\end{proof}

下一个引理反映了我们工作在单调情形这一事实。

\begin{lemma}
\label{lem:spectral-sequence-collapse}
谱序列在 $E_{N_L}$ 页坍缩，其中
\[
 N_L=\left[\frac{n+1}{\Sigma_L}\right].
\]
特别地，谱序列至多经过 $N_L$ 步就稳定。
\footnote{译者注：原书在本引理末句、随后关于 $E_\infty$ 的等式以及
定理 17.3.7 中均印作未定义的 $N$；本引理所定义并使用的是 $N_L$，
译文统一据此改正。}
\end{lemma}

\begin{proof}
这是~\eqref{eq:17-delta-prime-decomposition} 中给出的 $\delta$ 的结构的直接推论。
\end{proof}

这个引理使我们可以证明
\[
 E_\infty^{p,q}(\widehat C)=E_{N_L}^{p,q}(\widehat C).
\]

\begin{remark}
\label{rem:nonmonotone-projective-limit}
(FOOO09) 的一般性超出单调情形；对应谱序列并不在有限页坍缩，但其中证明了投射极限
\[
 \varprojlim E_r^{p,q}(\widehat C)
\]
存在，因而谱序列仍收敛。一般情形下该收敛性的证明见 (FOOO09)。
\end{remark}

现在令 $(\overline C,\overline\delta)$ 为 $(L,\phi_{\epsilon f}^1(L))$ 的局部\personFloer{}上同调复形；按定义，它同构于\personMorse{}复形 $(C^*(f),\delta_{\epsilon f})$。于是：

\begin{lemma}
\label{lem:E2-local-floer-morse}
存在同构
\[
 E_2^{*,*}(\widehat C)\cong
 H(\overline C;\overline\delta)\otimes_R
 \operatorname{gr}_*(F\Lambda_{(\omega,L)})
\]
作为 $\operatorname{gr}_*(F\Lambda_{(\omega,L)})$-模。
\end{lemma}

\begin{proof}
按定义，
\[
 E_1^{*,*}(\widehat C)\cong
 \overline C\otimes_R\operatorname{gr}_*(F\Lambda_{(\omega,L)}).
\]
由命题 17.1.3，$\delta_{(0)}=\overline\delta$。证明完毕。
\footnote{译者注：原书此处印作 $\delta_{(1)}=\overline\delta$；但局部\personFloer{}微分
此前记作 $\delta_{(0)}$，译文据此改正下标。}
\end{proof}

\begin{definition}
\label{def:filtration-on-homology}
定义 $F^qH(\widehat C,\delta)$ 为
$H(F^q\widehat C,\delta)$ 在 $H(\widehat C,\delta)$ 中的像。
\end{definition}

于是，由谱序列的一般结果，或由当前在有限页稳定的情形中的显式识别，得到：

\begin{theorem}
\label{thm:finite-page-associated-graded}
\[
 E_{N_L}^{p,q}(\widehat C)\cong
 F^qH^p(\widehat C,\delta)/F^{q+1}H^p(\widehat C,\delta)
\]
\footnote{译者注：原书分母漏印上标 $p$，因而不是分子所在分次部分的子模；
译文按同一定理的分子及定理 17.3.1(2) 补上 $p$。}
作为
$\Lambda^{(0)}(\lambda_0):=\Lambda_{(\omega,L)}^{(0)}
/F^{\lambda_0}\Lambda_{(\omega,L)}^{(0)}$-模。
\end{theorem}

按定义，
\[
 H^p(\widehat C,\delta)=
 HF^p(L,\phi_{\epsilon f}^1(L);\Lambda_{(\omega,L)}).
\]

\begin{remark}
\label{rem:novikov-residue-field}
回忆，在单调情形中 $\Lambda_{0,(\omega,L)}$ 是离散赋值环，其剩余域为
$\Lambda_{(\omega,L)}$。因此有同构
\begin{equation}
 \begin{aligned}
 HF(L,\phi_{\epsilon f}^1(L);\Lambda_{(\omega,L)})
 &\cong HF(L,\phi_{\epsilon f}^1(L);\Lambda_{0,(\omega,L)})\\
 &\qquad\otimes_{\Lambda_{0,(\omega,L)}}\Lambda_{(\omega,L)}.
 \end{aligned}
\label{eq:17-novikov-base-change}
\end{equation}
对子 $(L,\phi_{\epsilon f}^1(L))$，取子环 $\Lambda_{0,(\omega,L)}$ 上的\personFloer{}复形；即使在非单调情形中，也总能像对完整\personNovikov{}环
$\Lambda_{(\omega,L)}$ 那样构造它。但相伴同调
$HF(L,L';\Lambda_{0,(\omega,L)})$ 未必在哈密顿同痕下不变。另一方面，我们已经
证明，在以下情形中，单调对子 $(L,L')$ 的
$HF(L,L';\Lambda_{(\omega,L)})$ 在哈密顿同痕下不变：
\begin{enumerate}[label=(\arabic*)]
\item $L$ 与 $L'$ 的最小马斯洛夫数都大于 2（定理 16.4.6）；或
\item $L'$ 是 $L$ 的哈密顿形变，且 $L$ 的最小马斯洛夫数大于或等于 2
（定理 16.4.10）。
\end{enumerate}
\end{remark}

考察复形 $CF(L,\phi_{\epsilon f}^1(L))$ 的 $\delta$ 及这里构造的相伴谱序列，得到
作者在 (Oh96b) 中证明的如下定理。定理陈述中，在系数环 $\mathbb Z_2$ 上，我们把
$HF^*(L,\phi_{\epsilon f}^1(L))$ 简记为 $HF^*(L,L)$ 或
$HF(L;\Lambda_{(\omega,L)}^{\mathbb Z_2})$。

\begin{theorem}[与 (Oh96b) 比较]
\label{thm:large-minimal-maslov-HF}
设 $L\subset(M,\omega)$ 单调，且 $\Sigma_L\geq n+1$。考虑 $\mathbb Z_2$ 上的\personNovikov{}环 $\Lambda_{(\omega,L)}^{\mathbb Z_2}$ 及相伴\personFloer{}上同调
$HF^*(L,L)$。则：
\begin{enumerate}[label=(\arabic*)]
\item 若 $\Sigma_L>n+1$，则
$HF^*(L,L)\cong H^*(L;\mathbb Z_2)\otimes
\Lambda_{(\omega,L)}^{\mathbb Z_2}$；
\item 若 $\Sigma_L=n+1$，则对 $1<\ell<n$，
$HF^\ell(L,L)\cong H^\ell(L;\mathbb Z_2)$，并且或者
$HF^0(L,L)=HF^n(L,L)=\{0\}$，或者
\begin{equation}
 HF^n(L,L)=H^n(L;\mathbb Z_2),
 \qquad HF^0(L,L)=H^0(L;\mathbb Z_2).
\label{eq:17-end-degree-alternative}
\end{equation}
\end{enumerate}
\end{theorem}

\begin{proof}
因为 $n\geq1$，有 $\Sigma_L\geq2$，故 $HF(L,\phi(L))$ 良定义，并在 $\phi$ 的
哈密顿同痕下不变。以下简记为 $HF(L,L)$。

若 $\Sigma_L>n+1$，则 $N_L=0$，即 $\delta=\delta_{(0)}$。引理 17.3.3 与定理
17.3.7 蕴含
\[
 \begin{aligned}
 HF^*(L,L;\Lambda_{(\omega,L)}^{\mathbb Z_2})
 &\cong E_2^{*,*}(\widehat C)\\
 &\cong H(C(f);\delta_f)\otimes_{\mathbb Z_2}
 \operatorname{gr}_*(F\Lambda_{(\omega,L)}^{\mathbb Z_2})\\
 &\cong H^*(L;\mathbb Z_2)\otimes_{\mathbb Z_2}
 \Lambda_{(\omega,L)}^{\mathbb Z_2}
 \end{aligned}
\]
作为 $\Lambda_{(\omega,L)}^{\mathbb Z_2}$-模，故 (1) 得证。

若 $\Sigma_L=n+1$，则 $N_L=1$，故
$\delta=\delta_{(0)}+\delta_{(1)}\otimes q$，其中
$\delta_{(1)}:C^*(f)\to C^{*+n}(f)$。所以唯一非平凡贡献发生在 $C^0(f)$ 上。
在 $E_2$ 页，它诱导上边界映射
\begin{equation}
 \delta_{(1)}\otimes\Lambda_{(\omega,L)}^{(0);\mathbb Z_2}:
 H^0(f;\mathbb Z_2)\otimes\Lambda_{(\omega,L)}^{(0);\mathbb Z_2}
 \longrightarrow
 H^n(f;\mathbb Z_2)\otimes\Lambda_{(\omega,L)}^{(0);\mathbb Z_2},
\label{eq:17-E2-end-degree-map}
\end{equation}
并在其他所有次数上为零。因此 (2) 的第一个断言成立。对第二个断言，分
$\delta_{(1)}$ 是否为零两种情形。若不为零，它诱导同构，故
\[
 HF^0(L,L)=HF^n(L,L)=\{0\}.
\]
若为零，则~\eqref{eq:17-end-degree-alternative} 成立。对于（相对）spin 的 $L$，
上述定理可在 $\mathbb Z$ 系数上加强。
\end{proof}

最后简要提及 (FOOO09) 定义的\personFloer{}上同调上的乘积结构
\[
 \mathfrak m_2:HF(L,L;\Lambda_{(\omega,L)})\otimes
 HF(L,L;\Lambda_{(\omega,L)})\to HF(L,L;\Lambda_{(\omega,L)}).
\]
当 $\epsilon>0$ 充分小时，\personFloer{}上同调的\personMorse{}--\personBott{}版本记为
$HF(L,L;\Lambda_{(\omega,L)})$；可将它视为 $HF(L,\phi_{\epsilon f}^1(L))$。于是第
16.5 节构造的三角形乘积
\[
 \mathfrak m_2:HF(L,\phi_{\epsilon f}^1(L))\times
 HF(\phi_{\epsilon f}^1(L),L)\times HF(L,L)
\]
在令 $\epsilon\to0$ 后给出乘积
$\mathfrak m_2:\widehat C(L)\otimes\widehat C(L)\to\widehat C(L)$。由构造容易看出
\[
 \mathfrak m_2(F^{q_1}\widehat C(L)\otimes F^{q_2}\widehat C(L))
 \subseteq F^{q_1+q_2}\widehat C(L).
\]
这说明 $HF(L,L;\Lambda_{(\omega,L)})$ 上上述滤过由 $\mathfrak m_2$ 保持。下述
定理在 (FOOO09) 第 6.4.5 节中证明，证明见该处（另见 (Bu10)）。

\begin{theorem}
\label{thm:spectral-sequence-ring-compatible}
定理 17.3.1 构造的谱序列与环结构相容，具体如下。谱序列的每一页 $E_r$ 都有环结构
$\mathfrak m_2$，满足
\[
 \delta_r(\mathfrak m_2(x,y))
 =-\mathfrak m_2(\delta_r(x),y)
 +(-1)^{\deg x}\mathfrak m_2(x,\delta_r(y)).
\]
滤过 $F$ 与 $\mathfrak m_2$ 诱导的环结构相容。定理 17.3.1 中的同构都是环同构。
\end{theorem}

关于这个谱序列的进一步加强及其他应用，另见 (BCo09)、(D09)。
\stfhIndexRangeEnd{I179}{Puxulie}{谱序列（spectral sequence）}

\section{\personBiran{}与\personCieliebak{}定理}
\label{sec:biran-cieliebak-theorem}

从定义看，一个几何上显然的事实在\personFloer{}同调不变量的应用中起关键作用；这就是
下面的消失结论。

\begin{theorem}
\label{thm:displaceable-floer-vanishes}
假设 $L\subset(M,\omega)$ 带有一个在哈密顿同痕下不变的“\personFloer{}同调”
$HF(L;M)$。若 $L$ 可位移，即存在哈密顿微分同胚 $\phi$，使
$\phi(L)\cap L=\varnothing$，则 $HF(L;M)=\{0\}$。
\end{theorem}

\begin{proof}
任取满足 $L\cap\phi(L)=\varnothing$ 的哈密顿微分同胚 $\phi$。按
$CF(L,\phi(L))$ 的定义，$CF(L,\phi(L))=\{0\}$，故
$HF(L,\phi(L))=\{0\}$。另一方面，哈密顿同痕不变性给出
$HF(L,\phi(L))\cong HF(L;M)$。证明完毕。
\end{proof}

回忆，我们已在第 16 章证明：若下列条件之一成立，则\personFloer{}同调良定义并在哈密顿
同痕下不变：
\begin{enumerate}[label=(\arabic*)]
\item $L_0,L_1$ 满足 $\Sigma_{L_i}>2$；
\item $L_1$ 与 $L_0=L$ 哈密顿同痕，且 $L$ 满足 $\Sigma_L\geq2$。
\end{enumerate}
特别地，对情形 (2)：
\begin{enumerate}[label=(\arabic*)]
\item 对任意满足 $L\pitchfork\phi(L)$ 的哈密顿微分同胚 $\phi$，
$HF(L,\phi(L))$ 良定义；
\item 若给定二者之间的哈密顿同痕，则 $HF(L,\phi(L))$ 与
$HF(L,\phi'(L))$ 之间存在典范同构，而且该同构只依赖于这一（哈密顿）同痕类。
\end{enumerate}

\begin{corollary}
\label{cor:monotone-lagrangian-R2n-HF-zero}
对 $\mathbb R^{2n}$ 中任意满足 $\Sigma_L\geq2$ 的紧单调拉格朗日子流形，
$HF(L,\phi(L))$ 具有上述性质，因而 $HF(L,\phi(L))=\{0\}$。
\end{corollary}

\personBiran{}与\personCieliebak{}在 (BCi02) 中观察到，紧拉格朗日嵌入的类似自动可位移性质在任意
如下类型的辛流形上都成立：
\begin{equation}
 V\times M,
\label{eq:17-subcritical-product}
\end{equation}
其中 $V$ 是次临界\personWeinstein{}流形，$M$ 是任意紧辛流形。特别地，他们以定理
17.3.9 和推论 17.4.2 结合这一观察，用一种有趣方式研究了紧拉格朗日子流形
$L\subset V\times\mathbb{CP}^n$ 的拓扑性质。本节余下部分紧随 (BCi02) 的论证
叙述其定理。

先从 (EG91) 引入非紧辛流形研究中的几个标准概念。

\begin{definition}
\label{def:convex-at-infinity}
若辛流形 $(M,\omega)$ 带有一个在无穷远处为完备辛膨胀的向量场 $X$，则称它\termConvexAtInfinity{}。称向量场 $X$ 完备辛膨胀，是指 $X$ 的流
$\{\phi^t\}$ 完备，且满足 $(\phi^t)^*\omega=e^t\omega$。我们假设
$(M,\omega)$ 在无穷远处容许一个多重次调和穷竭函数。
\end{definition}

设 $(Q,\xi)$ 是接触流形，$\lambda$ 是 $\xi$ 的接触形式，即满足
$\ker\lambda=\xi$ 的 1-形式。与 $\lambda$ 相伴的\personReeb{}向量场 $X_\lambda$ 是唯一
满足
\begin{equation}
 X_\lambda\mathbin\lrcorner\lambda=1,
 \qquad X_\lambda\mathbin\lrcorner d\lambda=0
\label{eq:17-reeb-definition}
\end{equation}
的向量场。因此接触形式自然给出分裂
\[
 TQ=\mathbb R\{X_\lambda\}\oplus\xi.
\]
以 $\Pi=\Pi_\lambda:TQ\to TQ$ 表示满足
$\operatorname{Image}\Pi=\xi$、
$\ker\Pi=\mathbb R\{X_\lambda\}$ 的幂等算子，并以
$\pi=\pi_\lambda:TQ\to\xi$ 表示相应投影。

下面是接触流形研究中的标准定义。

\begin{definition}
\label{def:strong-symplectic-filling}
若存在紧辛流形 $(W,\omega)$，满足
\begin{enumerate}[label=(\alph*)]
\item $\partial W=Q$；
\item 在边界附近，$\omega=d\lambda$，其中 $\lambda$ 是某个 1-形式，并且
$\partial W$ 上 $\xi=\ker\lambda$；
\item $\lambda\wedge(d\lambda)^{n-1}$ 所定义的 $Q$ 的定向，与
$Q=\partial W$ 的边界定向一致，
\end{enumerate}
\stfhIndex{I182}{Qiangxinketianchong}{强辛可填充（strongly symplectically fillable）}
则称紧接触流形 $(Q,\xi)$ 强辛可填充，并称 $(W,\omega)$ 为 $(Q,\xi)$ 的\termStrongSymplecticFilling{}。
\end{definition}

根据辛邻域定理，可在 $W$ 中 $\partial W=Q$ 的领形邻域 $U_\delta$ 上选择函数
$r$，使得
\begin{equation}
 U_\delta\cong(1-\delta,1]\times Q,
 \qquad \omega=d(r\pi^*\lambda)\quad\text{于 }U_\delta,
\label{eq:17-collar-contact-form}
\end{equation}
其中投影 $\pi:M^{\geq R}\to Q$。再考虑柱体
$(1-\delta,\infty)\times Q$，沿带状区域
$(1-\delta,1]\times Q\cong U_\delta$ 作并
\begin{equation}
 \widehat W=W\#((1-\delta,\infty)\times Q).
\label{eq:17-symplectic-completion}
\end{equation}
把 $\omega$ 与 $(1-\delta,\infty)\times Q$ 上的 $d(r\pi^*\lambda)$ 粘合，就自然
延拓到 $\widehat W=:M$。按定义，每个无穷远处凸的辛流形都有这种分解；只需令
$\varphi=\log r$（等价地 $r=e^\varphi$），并使用等同
\[
 M^{\geq R}=M^{r\geq R}\cong[0,\infty)\times\varphi^{-1}(\ln R).
\]
用 $Q$ 上选定的接触形式 $\lambda$ 和投影 $\pi:M^{\geq R}\to Q$，辛形式写为
\begin{equation}
 \widehat\omega=d(r\pi^*\lambda)=d(e^s\pi^*\lambda).
\label{eq:17-cylindrical-symplectic-form}
\end{equation}
称 $(s,y)$ 为\termCylindricalCoordinates{}，称 $(r=e^s,y)$ 为\termConeCoordinates{}。

\begin{remark}
\label{rem:cone-coordinates-examples}
当
$M=\mathbb C^n\setminus\{0\}\cong(0,\infty)\times S^{2n-1}$ 时，
$(\sqrt r,y)$（$y\in S^{2n-1}$）就是 $\mathbb C^n\setminus\{0\}$ 的标准极坐标。
另一方面，若写
\[
 T^*N\setminus\{0\}\cong S^1(T^*N)\times\mathbb R_+,
\]
则对 $T^*N\setminus\{0\}\subset T^*N$ 的典范坐标 $(q,p)$，有 $r=|p|$。
\end{remark}

在柱体 $(-\infty,0]\times Q\subset(-\infty,\infty)\times Q$ 上有自然分裂
\[
 T_xM\cong\mathbb R\left\{\frac1r\frac\partial{\partial r}\right\}
 \oplus T_qQ
 \cong\operatorname{span}\left\{\frac\partial{\partial s},X_\lambda(q)\right\}
 \oplus\xi_q\cong\mathbb R^2\oplus\xi_q.
\]
以 $\widetilde X_\lambda$ 表示 $M$ 上唯一的向量场，它在平移下不变，与 $r$ 的
水平集相切，并投影到 $X_\lambda$。不致混淆时，也把它简记为 $X_\lambda$。

下面描述一族适配于接触型 $(Q,\xi)$ 流形 $M$ 的给定柱面结构的特殊殆复结构。先在
配备接触形式 $\lambda$ 的接触流形 $(Q,\xi)$ 上定义 CR-殆复结构。

\begin{definition}[接触三元组与三元组度量]
\label{def:contact-triad-metric}
设自同态 $J_Q:TQ\to TQ$ 满足 $J_Q^2=-\Pi$。若双线性形式
$d\lambda(\cdot,J_Q\cdot)$ 在 $\xi$ 上非退化，则称 $J_Q$ 为 $\lambda$-相容
CR-殆复结构；称三元组 $(Q,\lambda,J_Q)$ 为\termContactTriad{}。称 $Q$ 上由
\[
 g_Q(h,k)=g_{(\lambda,J_Q)}:=\lambda(h)\lambda(k)+d\lambda(h,J_Qk),
 \qquad h,k\in TQ,
\]
定义的度量为\termTriadMetric{}。
\end{definition}

\begin{definition}
\label{def:lambda-contact-type-J}
设 $(M,\omega)$ 是如上的带柱面端辛流形。若在端上，殆复结构 $J$ 分裂为
\[
 J=j\oplus J_Q|_\xi:T M\cong\mathbb R^2\oplus\xi
 \longrightarrow TM\cong\mathbb R^2\oplus\xi,
\]
其中 $J_Q|_\xi$ 与 $d\lambda|_\xi$ 相容，且
$j:\mathbb R^2\to\mathbb R^2$ 把 $\partial/\partial s$ 映到 $X_\lambda$，则称
$J$ 为\termLambdaContactType{}。
\end{definition}

在 $M$ 上选择关于驯服殆复结构 $J$ 的多重次调和穷竭函数 $\varphi$。还假设在一个
紧集之外，$J$ 在 $X$ 的流下不变。于是，对充分大的 $S$，水平集
$\varphi^{-1}(S)$ 带有诱导接触结构，并且流映射
\[
 \phi:(0,\infty)\times\varphi^{-1}(S)
 \to M\cap\varphi^{-1}([S,\infty)),
 \qquad(s,y)\longmapsto\phi^s(y)
\]
定义一个微分同胚。在这些坐标中 $\varphi=s$，且 $\phi^s$ 就是平移映射
$\phi^s(S,y)=(S+s,y)$。由此，无穷远处凸辛流形的定义也可借助接触流形
$N=\varphi^{-1}(S)$ 给出。

\begin{definition}
\label{def:weinstein-subcritical}
若非紧辛流形 $(V,\omega)$ 上存在穷竭\personMorse{}函数
$\varphi:V\to\mathbb R$，以及梯度型的完备\personLiouville{}向量场 $X$（即满足
\stfhIndex{I198}{Wensitanliuxing}{温斯坦流形（Weinstein manifold）}
$\mathcal L_X\omega=\omega$），则称 $V$ 为\termWeinsteinManifold{}。若
$\varphi$ 在每个临界点处的\personMorse{}指标都严格小于
$n=\frac12\dim V$，则称 $V$ 为\termSubcritical{}；否则称为临界的。
\end{definition}

\begin{exercise}
\label{ex:unstable-manifold-isotropic}
设 $(V,\omega,\varphi)$ 如定义 17.4.8。证明任意
$p\in\operatorname{Crit}\varphi$ 的不稳定流形都是迷向的。
\end{exercise}

\begin{corollary}
\label{cor:psh-morse-index-bound}
多重次调和函数任意临界点的\personMorse{}指标不大于
$n=\frac12\dim V$。
\end{corollary}

从次临界\personWeinstein{}流形的如下普遍性质开始。

\begin{proposition}[{(BCi02) 引理 3.2}]
\label{prop:compact-set-displaceable-subcritical}
设 $(V,\omega,X,\varphi)$ 是完备次临界\personWeinstein{}流形，$A\subset V$ 是任意紧
子集。则存在紧支哈密顿量 $H$，使得 $\phi_H^1(A)\cap A=\varnothing$。
\end{proposition}

\begin{proof}
以 $\Delta_X$ 表示 $X$ 的流的（正向）不变集。按定义，它是不稳定流形的并；这个并
是紧的，并且是维数不大于 $n-1$ 的有限个光滑流形之并。因此可找出支于
$\Delta_X$ 的一个小邻域中的哈密顿量 $K$，使
$\phi_K^1(\Delta_X)\cap\Delta_X=\varnothing$。由于 $\Delta_X$ 紧，可找出小邻域
$\Delta_X\subset U$，使 $\phi_K^1(U)\cap U=\varnothing$。

以 $\psi^t$ 表示 $X$ 的流。因为 $\mathcal L_X\omega=\omega$，$\psi^t$ 是共形辛的，
故共轭 $\psi^t\circ\phi_K^1\circ(\psi^t)^{-1}$ 是辛的。记
$h_t=\psi^t\circ\phi_K^1\circ(\psi^t)^{-1}$。固定充分大的 $T>0$，使
$\psi^{-T}(A)\subset U$。因为 $\phi_K^1$ 把 $U$ 从自身移开，得到
$h_T(A)\cap A=\varnothing$。最后验证 $h_T$ 由下列 $H=H(t,x)$ 生成：
\[
 H(t,x)=e^{tT}(\psi^{-tT})^*(K_t+G_t),
\]
其中 $G_t$ 是由方程（见命题 3.4.8 的证明）
\[
 dG_t=\lambda_V-(\phi_K^t)^*\lambda_V
\]
唯一确定的紧支哈密顿量，并可显式写作
\[
 G(t,x)=\int_0^t(K_s-dK_s(X))((\phi_K^s)^{-1}(x))\,ds.
\]
由此，因 $K$ 紧支，$H$ 也紧支。
\end{proof}

\begin{lemma}
\label{lem:product-lagrangian-displaceable}
设 $(V,\omega_V)$ 是满足 $\omega_V=d\lambda_V$ 的次临界\personWeinstein{}流形，
$(M,\omega_M)$ 是闭辛流形。则任意紧拉格朗日子流形
$L\subset V\times M$ 都可由紧支哈密顿同痕位移。
\end{lemma}

\begin{proof}
用平行于 $TV\oplus\{0\}$ 的哈密顿向量场位移 $L$。由于 $L$ 紧，它在 $V$ 上的
投影也紧，结论立即由命题 17.4.11 得到。
\end{proof}

\stfhIndex{I041}{Bilan-Qilibake}{比兰--齐利巴克（Biran--Cielibak）}
\personBiran{}与\personCieliebak{} (BCi02) 随后证明了如下同调球面性结论。

\begin{theorem}[\personBiran{}--\personCieliebak{}]
\label{thm:biran-cieliebak-spherical}
设 $V$ 是满足 $2\leq\dim V\leq2n$ 且 $c_1(V)=0$ 的连通次临界\personWeinstein{}流形。
考虑任意满足 $H_1(L,\mathbb Z)=\{0\}$ 的紧拉格朗日子流形
$L\subset V\times\mathbb{CP}^n$。则 $L$ 是\termCohomologicallySpherical{}，
即对所有 $1\leq\ell\leq\dim L-1$，
\[
 H^\ell(L;\mathbb Z_2)=\{0\}.
\]
\end{theorem}

\begin{proof}
首先，由 $c_1(V)=0$ 且 $V$ 是\personWeinstein{}（特别地，它正合），
$V\times\mathbb{CP}^n$ 与 $\mathbb{CP}^n$ 有相同单调常数。我们知道
$\mathbb{CP}^n$ 是单调辛流形，即在
$\pi_2(M)\cong H_2(M;\mathbb Z)$ 上
$I_\omega=\lambda_0I_{c_1}$，其中
$\lambda_0=\omega(L)/c_1(L)>0$。
\footnote{译者注：本式中的 $L$ 在本章一直表示待研究的拉格朗日子流形，不能作为
$\omega$ 与 $c_1$ 的二维同调类输入；原书没有在此明确命名应取的投影直线类，故
译文照录而不作不可唯一确定的修补。}

接着证明，任意满足 $H_1(L,\mathbb Z)=\{0\}$ 的紧拉格朗日嵌入
$L\subset V\times\mathbb{CP}^n=:M$ 都单调，且
$\Sigma_L\geq\dim V=2n+1>2$。由于 $H_1(L,\mathbb Z)=\{0\}$，自然映射
\[
 H_2(V\times\mathbb{CP}^n;\mathbb Z)
 \longrightarrow H_2(V\times\mathbb{CP}^n,L;\mathbb Z)
\]
满射，故 $H_2(V\times\mathbb{CP}^n,L;\mathbb Z)$ 中每个元素都来自
$H_2(V\times\mathbb{CP}^n;\mathbb Z)$。
\footnote{译者注：本定理假设 $2\leq\dim V\leq2n$，但原书证明这里写
$\Sigma_L\geq\dim V=2n+1$，随后又写
$\Sigma_L=2(n+1)=\dim L+1$；这些维数关系对一般允许的 $\dim V$ 并不一致。
译文照录并登记，不替原书改定理条件。}

现在考虑 $\pi_2(M,L)$ 上的同态 $I_{(\omega;L)},I_{(\mu;L)}$。二者都经由
$H_2(M,L;\mathbb Z)$ 分解，并在其中的挠元上为零。另一方面，模挠元有
\[
 H_2(V\times\mathbb{CP}^n;\mathbb Z)
 \cong H_2(V;\mathbb Z)\oplus H_2(\mathbb{CP}^n;\mathbb Z).
\]
所以只需对 $H_2(V;\mathbb Z)$ 与 $H_2(\mathbb{CP}^n;\mathbb Z)$ 的生成元证明
\begin{equation}
 I_{(\omega;L)}=\frac{\lambda_0}{2}I_{(\mu;L)}.
\label{eq:17-relative-monotonicity}
\end{equation}
对 $H_2(V;\mathbb Z)$，由于 $\omega$ 正合且假设 $c_1(V)=0$，这显然成立。

令 $[C_0]\subset\mathbb{CP}^n$ 为复直线的同调类
$[C_0]\in H_2(\mathbb{CP}^n;\mathbb Z)$，并把它视为 $H_2(M;\mathbb Z)$ 的元素；
以 $i_*[C_0]$ 表示其到 $H_2(M,L;\mathbb Z)$ 的推前。显然
\[
 I_{(\omega;L)}(i_*[C_0])=\omega([C_0])=I_\omega([C_0]),
\]
\[
 I_{(\mu;L)}(i_*[C_0])=2c_1([C_0])=2I_{c_1}([C_0]).
\]
回忆 $2c_1([C_0])=2(n+1)>2$。此外，
$\Sigma_L=2(n+1)=\dim L+1$，故 $L$ 满足定理 17.3.9 (2) 的假设。这蕴含
\[
 H^\ell(L;\Lambda_{(\omega,L)})
 \cong HF^\ell(L,L;\Lambda_{(\omega,L)}),
 \qquad1\leq\ell\leq\dim L-1.
\]
另一方面，因为 $V$ 次临界，$V\times\mathbb{CP}^n$ 上存在紧支哈密顿微分同胚
$\phi$，使 $L\cap\phi(L)=\varnothing$，即 $L$ 可位移。结合定理 17.3.9 (2)
与推论 17.4.2，即证定理 17.4.13。
\end{proof}

下一节将结合定理 17.4.1 与第 17.3 节构造的谱序列，导出各种推论。

\stfhIndex{I039}{Aodanwenti}{奥丹问题（Audin's question）}
\section{单调拉格朗日子流形的\personAudin{}问题}
\label{sec:audin-question-monotone}

除\personGromov{}的非正合性定理外，关于 $\mathbb C^n$ 中紧拉格朗日嵌入之辛拓扑，最早的
深刻问题之一是\personAudin{} (Au88) 提出的如下问题。

\begin{question}
\label{ques:audin-minimal-maslov-two}
任意紧拉格朗日环面的最小马斯洛夫数是否都是 2？
\end{question}

\stfhIndex{I154}{Boerjieluoweiqi}{波尔捷罗维奇（Polterovich）}
\stfhIndex{I193}{Weitaibo}{维泰博（Viterbo）}
第一个肯定答案由\personViterbo{} (Vi90) 在一般情形下得到，而\personPolterovich{} (Po91a) 则在
$\mathbb R^4$ 中得到。\personViterbo{}使用作用量泛函的经典临界点理论，\personPolterovich{}使用
伪全纯曲线方法。对单调拉格朗日环面 $L\subset\mathbb C^n$，作者利用拉格朗日\personFloer{}同调的良定义性及厚--薄二分思想，在 $n\leq24$ 时肯定地回答了这个问题
(Oh96b)。同一定理也适用于一般 $(M,\omega)$ 中任意可位移单调拉格朗日环面。

\personFukaya{}、\personOh{}、\personOhta{}、\personOno{} (FOOO09) 与 L. \personBuhovsky{} (Bu10) 独立地利用谱序列与
$\mathfrak m_2$ 诱导乘积的相容性，去掉了维数限制 $n\leq24$。使用谱序列乘法性质
的想法来自 P. \personSeidel{}（私人通信，2002）。结果如下；这里的证明借自 (FOOO09)。

\begin{theorem}
\label{thm:audin-monotone-displaceable-torus}
设 $L\subset(M,\omega)$ 是可位移单调拉格朗日环面。则其最小马斯洛夫数满足
\[
 \Sigma_L=2.
\]
\end{theorem}

\begin{proof}
先指出 $T^n$ 可定向且为 spin，而单调性依定义蕴含 $\Sigma_L>0$。因为 $T^n$
可定向，非零整数 $\Sigma_L$ 必为偶数，故 $\Sigma_L\geq2$。

用反证法。反设 $\Sigma_L\geq3$。（由于 $T^n$ 可定向，这其实蕴含
$\Sigma_L\geq4$，但证明不用此事实。）单调性于是蕴含
$HF(L,L;\Lambda_{(\omega,L)})$ 良定义并在哈密顿同痕下不变。

边界映射分解为
\[
 \delta=\delta_{(0)}+\delta_{(1)}+\cdots+\delta_{(N)}
 =:\delta_{(0)}+\delta',
\]
其中 $N\leq[(n+1)/\Sigma_L]$，而 $\delta_{(k)}$ 具有形式
\begin{equation}
 \delta_{(k)}=\overline\delta_{(k)}\otimes T^{k\lambda_0}e^{k\Sigma_L}.
\label{eq:17-delta-k-monotone-torus}
\end{equation}
这里 $\overline\delta_{(k)}:E_k\to E_k$ 的次数为 $1-k\Sigma_L$。由于
$\Sigma_L\geq3$，当 $k\geq1$ 时该次数不大于 $-2$。

现在考虑 $\delta$ 在 $E_2$ 项
\[
 E_2=H^*(T^n;\mathbb Q)\otimes\Lambda_{0,\mathrm{Nov}}
\]
上的作用。回忆，上同调环 $H^*(T^n;\mathbb Q)$ 由一次上同调类
$\alpha_1,\ldots,\alpha_n\in H^1(L;\mathbb Q)$ 生成。简单的次数计算给出
\[
 \delta_2(\alpha_i)=\sum_{k=1}^{N}\overline\delta_{(k)}(\alpha_i)
 \cdot T^{k\lambda_0}e^{k\Sigma_L}=0,
 \qquad i=1,\ldots,n.
\]
\footnote{译者注：原书在本式中用同一个 $k$ 同时编号上同调生成元
$\alpha_k$ 与微分分量 $\overline\delta_{(k)}$，只写出一个匹配下标的分量；
但 $\delta_2$ 是所有正能量分量之和。译文改用 $i$ 编号生成元并补出对
$1\leq k\leq N$ 的求和；每一分量因次数至多为 $-2$ 都在 $H^1$ 上消失。}
因此
\[
 \delta_2(\mathfrak m_2(\alpha_i,\alpha_j))
 =-\mathfrak m_2(\delta_2(\alpha_i),\alpha_j)
 -\mathfrak m_2(\alpha_i,\delta_2(\alpha_j))=0.
\]
另一方面，由能量考察，
\[
 \mathfrak m_2(\alpha_i,\alpha_j)\equiv
 \alpha_i\smile\alpha_j\pmod{T^{\lambda_0}}.
\]
故对所有 $i,j=1,\ldots,n$，都有
$\delta_2(\alpha_i\smile\alpha_j)=0$。对 $\alpha_i$ 的所有可能乘积归纳应用上述
论证，得到谱序列在 $E_2$ 项退化，从而
\[
 HF^*(L;\Lambda_{0,\mathrm{Nov}})
 \cong H^*(T^n;\mathbb Q)\otimes\Lambda_{0,\mathrm{Nov}},
\]
特别地，它是 $\Lambda_{0,\mathrm{Nov}}$ 上的自由模。因此
\begin{equation}
 HF^*(L;\Lambda_{\mathrm{Nov}})
 \cong H^*(T^n;\mathbb Q)\otimes\Lambda_{\mathrm{Nov}}\neq0.
\label{eq:17-torus-floer-nonzero}
\end{equation}
（到目前为止尚未使用 $L\subset M$ 可位移的假设；它将在证明的最后一步使用。）
另一方面，由 $HF^*(L;\Lambda_{\mathrm{Nov}})$ 的不变性及存在满足
$\phi(L)\cap L=\varnothing$ 的哈密顿微分同胚 $\phi$，可知
$HF^*(L;\Lambda_{\mathrm{Nov}})=\{0\}$。这与
\eqref{eq:17-torus-floer-nonzero} 矛盾，因而 $\Sigma_L\leq2$ 得证。
\end{proof}

\begin{remark}
\label{rem:audin-modern-perspective}
从\personFloer{}理论的现代观点看，\personViterbo{}在 (Vi90) 中的证明可以重述为辛流形环路空间上
哈密顿\personFloer{}同调的一个应用 (Vi99)，而\personPolterovich{}的证明可视为拉格朗日\personFloer{}同调的应用。长期以来，人们颇感困惑：为何这两种\personFloer{}同调会给出同类结果？后来
\stfhIndex{I093}{Shenguxianzhi}{深谷贤治（Fukaya）}
\personFukaya{} (Fu06) 在拉格朗日子流形的环路空间上引入\personChas{}--\personSullivan{}括号，并把
$L\subset(M,\omega)$ 的 $A_\infty$-代数重构为德拉姆复形之 DGA 的形变，从而
在完全一般的情形下肯定地回答了\personAudin{}问题。此外还有 Y. \personEliashberg{}倡导的辛场论
方法（私人通信，2001）。据作者所知，这些证明的细节尚未发表。
\end{remark}

\section{\texorpdfstring{$\Ham(S^2)$}{Ham(S2)} 上的\personPolterovich{}定理}
\label{sec:polterovich-Ham-S2}

联系\personHofer{}距离，一个自然问题是：对给定紧辛流形，$\Ham(M,\omega)$ 的直径究竟
有限还是无限？

\begin{definition}[\personHofer{}直径]
\label{def:hofer-diameter}
设 $(M,\omega)$ 是辛流形，$\Ham(M,\omega)$ 是其哈密顿微分同胚群。（若
$(M,\omega)$ 是开的，则改取 $\Ham_c(M,\omega)$。）定义
\[
 \operatorname{diam}(\Ham(M,\omega))
 =\sup_{\phi\in\Ham(M,\omega)}\|\phi\|.
\]
\end{definition}

\stfhIndex{I154}{Boerjieluoweiqi}{波尔捷罗维奇（Polterovich）}
\personOstrover{} (Os03) 与\personPolterovich{} (Po98b) 分别证明：在
$\pi_2(M)=\{0\}$ 的非球面情形以及 $S^2$ 的情形，\personHofer{}直径为无穷。

\begin{theorem}[\personOstrover{}]
\label{thm:ostrover-hofer-infinite}
设 $(M,\omega)$ 是满足 $\pi_2(M)=0$ 的闭辛流形。则
\stfhIndex{I111}{Huofeierzhijing}{霍费尔直径（Hofer's diameter）}
$\Ham(M,\omega)$ 的\termHoferDiameter{}为无穷。
\end{theorem}

\personUsher{}在 (Ush10a) 中给出了这一结果的各种推广。另一方面，\personPolterovich{}证明了：

\begin{theorem}[\personPolterovich{}]
\label{thm:polterovich-hofer-S2-infinite}
$\Ham(S^2)$ 的\termHoferDiameter{}为无穷。
\end{theorem}

\personOstrover{}的定理使用哈密顿\personFloer{}同调证明，而\personPolterovich{}使用拉格朗日交\personFloer{}同调。我们将在第四部分回到\personOstrover{}的证明。本节紧随\personPolterovich{}的原始论证，
给出他的证明。

先从应用于哈密顿环路的拉格朗日悬垂构造开始。考虑满足
\[
 H:[0,1]\times M\to\mathbb R,
 \qquad \phi_H^1=\mathrm{id}
\]
的哈密顿量。考虑任意闭辛流形 $(M,\omega)$。设
$\phi\in\Ham(M,\omega)$，$F$ 是满足 $F\mapsto\phi$ 的归一化哈密顿量。

\begin{definition}
\label{def:periodic-hamiltonian-loops}
以 $\mathcal H_{\mathrm{loop}}$ 表示生成哈密顿环路的一周期哈密顿量 $H$（即
$H(t,x)=H(t+1,x)$）的集合；换言之，$\phi_H$ 满足
$\phi_H^1=\mathrm{id}$。
\end{definition}

回忆，对任意给定的哈密顿环路，总能重新参数化它，使
$\phi_H^t\equiv\mathrm{id}$ 在 $t=0,1$ 附近成立，从而使相伴哈密顿量 $H$ 为
一周期的。特别地，对某个常数 $c$，在
$[0,\delta]\cup[1-\delta,1]\times M$ 上 $H\equiv c$。

对定义 3.7.6 中哈密顿同痕 $\phi_H$，考虑赤道
$L\subset(M,\omega)$ 的悬垂，记作
\[
 N(L,\phi_H)\subset T^*L\times M,
\]
并令 $N_0=N(L,\mathrm{id})$。注意 $N_0$ 就是积 $o_L\times L$，且命题 3.7.7
蕴含 $\mathcal L=N(L,\phi_H)$ 与 $N_0$ 哈密顿同痕。此外，因为
$L\subset(M,\omega)$ 单调且最小马斯洛夫数为 2，$N_0$ 亦然。所以
$N_0$ 的\personFloer{}同调 $HF(N_0,N(L,\phi_H))$ 在\personNovikov{}环
$\Lambda^{\mathbb Z_2}$ 上良定义，并由定理 16.4.10 同构于 $HF(N_0,N_0)$。
另一方面，我们有如下计算。

\begin{lemma}
\label{lem:floer-product-suspension}
假设 $L\subset(M,\omega)$ 单调且 $\Sigma_L\geq2$。则
\[
 HF(N_0,N_0)\cong HF(o_L,o_L)\otimes HF(L,L;M).
\]
\end{lemma}

\begin{proof}
由于 $N_0$ 单调且最小马斯洛夫数至少为 2，$N_0$ 与其哈密顿形变的相伴\personFloer{}同调
不依赖于同痕和殆复结构的选择（再用定理 16.4.10）。
\footnote{译者注：本引理只假设 $\Sigma_L\geq2$，而原书断言积拉格朗日
$N_0=o_L\times L$ 的最小马斯洛夫数恒为 2；一般只能由积结构得到至少为 2。
译文按证明实际所需的下界改正。}

注意拉格朗日子流形 $N_0=o_L\times L$ 是积辛流形 $T^*L\times M$ 中的直积。
利用 $HF$ 在哈密顿同痕下的不变性（仍用定理 16.4.10），可在
$T^*L\times M$ 上考虑积殆复结构 $J_0$，以及 $N_0$ 的通有积哈密顿同痕
$\phi_{H'}=\phi_{H_1'}\times\phi_{H_2'}$。相伴\personFloer{}流线方程分裂，因此相应\personFloer{}复形 $(CF(N_0,\phi_{H'}^1(N_0)),\partial)$ 成为
\[
 CF(N_0,\phi_{H'}^1(N_0))=CF(o_L,\phi_{H_1'}^1(o_L))
 \otimes CF(L,\phi_{H_2'}^1(L)),
\]
\footnote{译者注：原书在本式左端突然写未定义的 $L'$；由上一句正在讨论的\personFloer{}复形可唯一确定其应为 $\phi_{H'}^1(N_0)$，译文据此补正。}
其边界映射为
\[
 \partial(x\otimes y)=\partial(x)\otimes y+x\otimes\partial(y).
\]
于是立即得到同构
\[
 HF(N_0,N_0)\cong HF(o_L,o_L)\otimes HF(L,L;M).
\]
\footnote{译者注：引理陈述给出张量积 $\otimes$，且证明中的链复形也是张量积，
但原书证明末句却印作直和 $\oplus$；译文据引理陈述及前一公式改正。}
\end{proof}

\stfhIndex{I181}{Wendingxiangjiaoxingzhi}{稳定相交性质（stable intersection property）}
直接推论是：若 $L\subset M$ 满足 $HF(L,L;M)\neq0$，则它具有如下稳定相交性质。

\begin{corollary}
\label{cor:stable-intersection-from-HF}
假设 $L\subset(M,\omega)$ 单调且 $\Sigma_L\geq2$，因而\personFloer{}同调
$HF(L,L;M)$ 有定义。再假设 $HF(L,L;M)\neq0$。则对 $(M,\omega)$ 的任意哈密顿
环路 $\phi_H$，都有
\[
 (o_L\times L)\cap N(L,\phi_H)\neq\varnothing.
\]
\end{corollary}

从下面的定义看，这样的拉格朗日子流形 $L$ 具有稳定拉格朗日交性质。

\begin{definition}
\label{def:stable-lagrangian-intersection}
若对任意哈密顿环路 $\phi_H$（即满足 $\phi_H^1=\mathrm{id}$），都有
$(o_L\times L)\cap N(L,\phi_H)\neq\varnothing$，则称拉格朗日子流形
$L\subset(M,\omega)$ 具有\termStableLagrangianIntersectionProperty{}。
\end{definition}

如下\personFloer{}同调非消失结果实质上已在 (Oh93b) 中证明。

\begin{theorem}
\label{thm:equator-floer-nonzero}
考虑 $S_{\mathrm{eq}}^1\subset S^2$。则
\[
 HF(S_{\mathrm{eq}}^1,S_{\mathrm{eq}}^1;S^2)
 \cong H(S^1,\Lambda^{\mathbb Z_2}).
\]
特别地，$S_{\mathrm{eq}}^1$ 具有稳定拉格朗日交性质。
\end{theorem}

\begin{proof}
在 $S_{\mathrm{eq}}^1$ 上取一个 $C^2$-小\personMorse{}函数 $f$；它恰在赤道的一对对跖点
有两个临界点，并生成给定赤道的一个小旋转。
\footnote{译者注：原书把这两个位于赤道上的临界点称为 $S^2$ 的南极和北极，但
南北极不在 $S_{\mathrm{eq}}^1$ 上；证明只需赤道上的一对对跖临界点，译文据此改正。}
应用\personRiemann{}映射定理可直接检验，在本情形
中厚--薄分解给出
\[
 \partial=\partial_f^{\mathrm{Morse}}+\partial',
\]
其中 $\partial'$ 由两条\personFloer{}流线贡献：从上下半球切掉
$S_{\mathrm{eq}}^1$ 与 $\phi_f^1(S_{\mathrm{eq}}^1)$ 之间的“薄”流线，就得到它们。
恰有两条这样的“厚”流线。因此矩阵系数为
\[
 T^{\omega(u_1)}+T^{\omega(u_2)},
\]
其中 $\omega(u_i)$ 是相应\personFloer{}流线 $u_i$ 的面积。另一方面，由
$S_{\mathrm{eq}}^1$ 的单调性，两个面积相同，故矩阵系数变为 $2T^\lambda$，其中
$\lambda$ 是公共面积。因此在 $\Lambda^{\mathbb Z_2}$ 中该矩阵系数为零。这说明
$\partial=\partial_f^{\mathrm{Morse}}$，证明完毕。
\end{proof}

接下来回忆拉格朗日悬垂的显式定义
\[
 \iota_{\phi_H,L}:(t,x)\longmapsto
 (t,-H(t,\phi_H^t(x)),\phi_H^t(x)),
\]
其像就是 $N(L,\phi_H)$。交点
\[
 (t,a,y)\in N_0\cap N(L,\phi_H)
\]
由方程
\begin{equation}
 0=-H(t,y)=a,
 \qquad y\in\operatorname{Fix}\phi_H^t\cap L
\label{eq:17-suspension-intersection-characterization}
\end{equation}
刻画。

还回忆，在 $[-\delta,\delta]\times M$ 上 $H\equiv c$，其中把
$[-\delta,\delta]$ 看作 $S^1=\mathbb R/\mathbb Z$ 的子集。为去掉平凡交点，选取
可任意小的 $c\neq0$，再重新归一化 $H_t$，使其满足归一化条件。做到这一点时可让
$\|H\|$ 任意小。

有了这些准备，可以完成定理 17.6.3 的证明。

\begin{proof}[定理 17.6.3 的证明]
先用周期哈密顿量集合 $\mathcal H_\phi$（见定义 17.6.4）描述
$\phi\in\Ham(M,\omega)$ 的\personHofer{}范数 $\|\phi\|$。

\begin{lemma}
\label{lem:hofer-norm-loop-infimum}
设 $\phi\in\Ham(M,\omega)$ 且 $\phi=\phi_F^1$。则
\[
 \|\phi\|=\inf_{H\in\mathcal H_{\mathrm{loop}}}\|F-H\|.
\]
\end{lemma}

\begin{proof}
证明留作习题。
\end{proof}

考虑哈密顿同痕 $\phi_H$ 所给赤道
$L=S_{\mathrm{eq}}^1\subset S^2$ 的悬垂，记为
\[
 N(S_{\mathrm{eq}}^1,\phi_H)\subset T^*S^1\times S^2,
\]
其中
\[
 S^1=S^1(2)=\mathbb R/2\mathbb Z,
\]
并令 $N_0=N(S_{\mathrm{eq}}^1,\mathrm{id})$。
\footnote{译者注：原书此式在商号后印有额外的 $c$，写作
$\mathbb R/c\,2\mathbb Z$；此前的 $c$ 是可任意小的哈密顿常数，不可能定义固定的
$S^1(2)$，故译文改为 $\mathbb R/2\mathbb Z$。}
注意 $N_0$ 就是积 $o_{T^*S_{\mathrm{eq}}^1}\times S_{\mathrm{eq}}^1$；命题 3.7.7
蕴含，只要哈密顿环路在 $\Ham(S^2,\omega)$ 中可缩，
$\mathcal L=N(S_{\mathrm{eq}}^1,\phi_H)$ 就与 $N_0$ 哈密顿同痕。

另一方面，$\pi_1(\Ham(S^2))\cong\mathbb Z_2$，且沿垂直于给定赤道
$S_{\mathrm{eq}}^1$ 的轴作整周旋转生成一个非平凡环路。这个环路把
$S_{\mathrm{eq}}^1$ 映到自身，且相伴哈密顿量在 $S_{\mathrm{eq}}^1$ 上消失。因此
相伴悬垂仍为 $N_0$。所以，对所有哈密顿环路 $\phi_H$，
$N(S_{\mathrm{eq}}^1,\phi_H)$ 都与 $N_0$ 哈密顿同痕。

又因 $S_{\mathrm{eq}}^1\subset S^2$ 单调且最小马斯洛夫数为 2，$N_0$ 亦然。因此
$N_0$ 的\personFloer{}同调
\[
 HF(N_0,N(S^1,\phi_H))
\]
在\personNovikov{}环 $\Lambda^{\mathbb Z_2}$ 上良定义，并由定理 16.4.11 同构于
$HF(N_0,N_0)$。
然后对
\[
 HF(o_{T^*S^1},o_{T^*S^1})\cong H(S^1)
 \cong\Lambda^{\mathbb Z_2}\oplus\Lambda^{\mathbb Z_2}
\]
应用余切丛同构，并对
\[
 HF(S_{\mathrm{eq}}^1,S_{\mathrm{eq}}^1;S^2)
 \cong H(S_{\mathrm{eq}}^1)
 \cong\Lambda^{\mathbb Z_2}\oplus\Lambda^{\mathbb Z_2}\neq0
\]
应用定理 17.6.8。由推论 17.6.6，$S_{\mathrm{eq}}^1$ 具有稳定拉格朗日交性质。

于是由~\eqref{eq:17-suspension-intersection-characterization}，存在
$(t,y)\in S^1\times L$，使得 $H(t,y)=0$。注意上述证明依赖于
$\Ham(S^2)$ 的基本群有限这一事实。它与引理 17.6.9 立即给出如下推论。

\begin{corollary}
\label{cor:hofer-lower-bound-equator}
设归一化哈密顿量 $F$ 满足
$F|_{S_{\mathrm{eq}}^1}\equiv c>0$，并令 $\phi=\phi_F^1$。则
$\|\phi\|\geq c$。
\end{corollary}

最后，在此推论中可选择归一化哈密顿量 $F$，使 $c$ 任意大，证明即告完成。
\end{proof}

事实上，只要任意闭 $(M,\omega)$ 包含一个具有稳定相交性质的紧拉格朗日子流形，
上述证明就适用。

\begin{theorem}
\label{thm:stable-intersection-implies-infinite-\personHofer{}}
假设 $(M,\omega)$ 包含一个紧拉格朗日子流形 $L$，并且 $L$ 具有\termStableLagrangianIntersectionProperty{}。则
$\operatorname{diam}\Ham(M,\omega)$ 为无穷。
\end{theorem}

在这方面，我们已经在上述证明中说明
$S_{\mathrm{eq}}^1\subset S^2$ 具有稳定拉格朗日交性质。用同一论证还可推出，
$\mathbb{RP}^n$ 以及 $\mathbb{CP}^n$ 中的\personClifford{}环面 $T^n$ 也具有这一性质。

\begin{remark}
\label{rem:hofer-versus-spectral-diameter}
定理 17.6.3 与将在第 22.3 节介绍的谱直径有限性定理形成鲜明对比。我们将在第
22.5 节证明
\[
 \operatorname{diam}_\rho(\Ham(S^2))\leq4\pi
 \quad(=S^2\text{ 的面积}).
\]
\end{remark}

